\section{The current arithmetic route and its next calculation}
\label{sec:research-state}
The source chain now supplies an actual \(\mathbb C[t]\)-module
cohomology \(Q\), a finite invariant packet inside that cohomology,
a specified nonzero extension class, an original theta boundary
for its representative defect, and an exact arithmetic formula
for the relative weight form. The local continuation makes that
last formula an explicit confluent interpolation problem.
These constructions are logically prior to any purity conclusion.
The subsequent continuation now supplies certified entries of the
original arithmetic moment matrix, the exact finite-compression
and infinite-tail resolvent comparison, and a convergence rate for
the initial resolvent including its complete off-line jets.

\subsection{A single chain of exact comparison maps}
For each actual finite packet, retaining every full order, the
maps proved in this reader are
\[
 E=\mathbb C[t]/h
 \xrightarrow{\,R_{\mathrm{ref}}\,}\mathscr B
 \xrightarrow{\,q\,}Q,\qquad
 J_ZR_{\mathrm{ref}}=1,\qquad
 DR_{\mathrm{ref}}-R_{\mathrm{ref}}A
       =\Theta\phi_*\ell .
\]
Every admitted polynomial correction keeps these same quotient
and jet maps. At its distinguished orthogonal level let \(B_{N,u}(t)=(p_{N,u}(t)-R_Z(u)(t))/h(t)\), a polynomial of degree at most \(m\). The complete calculation is
\[
\begin{gathered}
 d\nu_Z(t)=\left|\frac{g(1/2+it)}{h(1/2+it)}\right|^2
       \frac{dt}{2\pi},\quad N=d+m,\quad
 C_NP=j_hP,\quad
 K_N=C_N(H_N^\nu)^{-1}C_N^*,\\
 p_{N,u}=(H_N^\nu)^{-1}C_N^*K_N^{-1}U_\varepsilon u,\qquad
 R_mu=R_{\mathrm{ref}}u+
       \Theta(B_{N,u}(D)\phi_*),\\
 G_m=U_\varepsilon^*K_N^{-1}U_\varepsilon,\qquad
 W_m=U_\varepsilon^*K_N^{-1}
     (AK_N+K_NA^*-K_N)K_N^{-1}U_\varepsilon .
\end{gathered}
\tag{R1}
\]
The middle line evaluates the specified polynomial at the original \(D\); every source vector lies in \(V\).

Write \(v_n=j_hq_n\) and \(\kappa_n=\|q_n\|_{\nu_Z}^2\).
The original arithmetic measure now determines the exact quantity
\[
 \epsilon_m=
 \frac{|\operatorname{Re}(v_N^*K_N^{-1}v_{N+1})|
 +\sqrt{(v_{N+1}^*K_N^{-1}v_{N+1})(v_N^*K_N^{-1}v_N)
       -(\operatorname{Im}(v_N^*K_N^{-1}v_{N+1}))^2}}
      {\kappa_N}.
 \tag{R2}
\]
The proofs in Section~\ref{sec:tau-boundary-kernel} establish
(R1)--(R2) before passing to cohomology.
For a reflection-stable packet the first absolute-value term is
zero by the exact trace identity. The same sequence obeys
\[
 \epsilon_m\ \geq\
       \sum_{\rho\in Z}m_\rho
                   |\operatorname{Re}\rho-\tfrac12|.
 \tag{R3}
\]
Thus the rank-two boundary detects the combined displacement of
the entire packet, not only its largest individual displacement.
The estimate concerns the distinguished \(R_m\) sequence.
The larger family of all polynomial source corrections has the
separately proved metric image and sharp infimum in
Section~\ref{sec:tau-chain-new}; those results do not identify
its arbitrary member with \(R_m\).

The full two-row source observation in
Section~\ref{sec:tau-boundary-kernel} is injective after the
original \(A\)-iterates are retained. It therefore offers a
finite arithmetic calculation preserving every nilpotent jet:
compute the source moments, the adjacent monic polynomials and
their complete residue vectors, rather than discard a component
whose scalar trace vanishes.

\subsection{Topology, duality and the two sides of a weight estimate}
The actual source moments now have the absolutely convergent
incomplete-gamma expression and explicit double-tail bound of
Section~\ref{sec:arithmetic-moments}. Its certified six-by-six
Gram matrix concerns the original \(|2\xi|^2dt/(2\pi)\) measure.
The exact identity (AM.19) supplies the original monic recurrence
norms through degree five, with interval enclosures.

For the \(|g/h|^2\) kernel, Section~\ref{sec:kernel-resolvent}
proves a universal polynomial quadrature formula, including
centres that coincide with a quadrature node. On the precise
nonintersection domain its alternate expression is
\[
 K_{n-1}=\kappa_{n-1}^{-1}q_n(A)S_nq_n(A)^*,\qquad
 AS_n+S_nA^*-S_n=cr_n^*+r_nc^*,\quad r_n=F_n(A)c.
\]
The exact Schur comparison (KR.15)--(KR.20) then joins the initial
and terminal observations to the actual infinite-dimensional tail.
Section~\ref{sec:gaussian-resolvent-convergence} proves the explicit
initial-resolvent rate (GC.8) and its complete finite-algebra
bound (GC.23). The projection to that off-line domain retains
the critical summand as its stated kernel. All factors in the
Schur identity remain in converting the initial estimate to the
terminal form that determines \(\epsilon_m\).

In particular, Section~\ref{sec:metric-departure} identifies
exactly the two contributions in a reflection-stable packet:
\[
 \epsilon_m^2=2\sum_{\rho\in Z}m_\rho
       (\operatorname{Re}\rho-1/2)^2+\mathfrak d_{G_m}(A).
\]
Its full-jet inequalities also quantify the required metric
ratios along each retained nilpotent chain. These identities
do not infer a relative estimate from absolute matrix shrinkage.

The theta-polynomial span is dense in every finite differential
graph norm. The completed jet graph, however, is
\(\mathscr H^{(r)}\oplus E\), with quotient \(E\), not a
quotient of \(\mathscr H^{(r)}\) in which the arithmetic
component has vanished. This supplies the exact comparison
between the shrinking Hilbert representative and its unchanged
cohomological class.

The exponential observation norm supplies a second exact
comparison. A Mellin jet at \(\rho\) is bounded precisely in
the strip
\[
 |\operatorname{Re}\rho-\tfrac12|<a,
\]
with the factorial norm in Theorem~\ref{thm:weighted-jet}.
The original differential then has the energy term
\(2a\,\operatorname{sgn}(\log x)\) in
Theorem~\ref{thm:weighted-energy}.
Consequently stronger topology and weight control are related
by that explicit term; choosing a positive observation norm
does not remove it.

For any packet \(Z\), reflection gives the conjugate-linear ring
isomorphism
\[
 \dagger_Z:E_Z\longrightarrow E_{Z^\dagger},\qquad
 f^\dagger(s)=\overline{f(1-\overline s)},\quad
 Z^\dagger=\{1-\overline\rho:\rho\in Z\},
\]
with all multiplicities retained. In fact
\(h_Z^\dagger=(-1)^dh_{Z^\dagger}\); this proves well-definedness,
and applying the same operation twice is the identity. Restrict
the following internal residue pairing to reflection-stable
\(Z\), so that this map is an involution on \(E_Z\).
The arithmetic residue form is the source form
\[
 \mathcal R_Z(f,u)=\sum_{\rho\in Z}
             \operatorname{Res}_\rho\frac{f^\dagger u}{g}\,ds .
\]
Its Jacobian contraction is
\[
 \mathcal R_Z(f,g'u)
 =\sum_{\rho\in Z}m_\rho\overline{f(1-\overline\rho)}u(\rho).
 \tag{R4}
\]
Indeed locally \(g=z^{m_\rho}w(z)\) gives
\(g'/g=m_\rho/z+w'/w\), proving (R4) with multiplicities.
The trace form pairs reflected centres and can retain
nontrivial null directions from nilpotents.
In the retained basis \(1,t,\ldots,t^{d-1}\), define
\[
 S_{ij}=\sum_{\rho\in Z}\operatorname{Res}_\rho
       \frac{(s^i)^\dagger s^j}{g(s)}\,ds,
 \quad 0\leq i,j<d,\qquad J_g=M_{j_hg'}.
\]
Changing representatives by \(h\) makes the local quotient by
\(g\) holomorphic, proving that this matrix represents the
intrinsic form. By (R4) the trace pairing has matrix \(SJ_g\).
Its exact comparison to the positive source metric is therefore
\(G_m^{-1}SJ_g\), since
\(f^*G_m(G_m^{-1}SJ_g)u=f^*SJ_gu\).
Its positivity is an arithmetic question about this operator,
not a consequence of the positive Gram construction.

The finite-field two-sided weight theorem remains a source for
methods, with its compact-support comparison, duality and
tensor hypotheses. This reader has not constructed a realization
meeting those geometric hypotheses for the entire arithmetic
quotient. It has instead calculated the current characteristic-zero
objects and their original maps, including the local-cohomology
spectral extraction and meromorphic complement documented in
the complete finite-comparison appendix.

\subsection{Provenance and mathematical status}
The coherent interpolation now has the full coefficient identity
\(J_{h,N}=U_{\upsilon_h}C_N\), with
\(\upsilon_h=j_h(g/h)\) and \(\varepsilon_h=\upsilon_h^{-1}\).
Equations (CT.3)--(CT.7) prove equality with (R1), including the
original representative, and compatibility with the complete CRT
injections between finite packets.  The actual reflection preserves
the unique minimum and proves the same-metric signed pairing of
(CT.10) and (CT.16a).

On the original tensor complex, the complete calculation is
\[
 \begin{gathered}
 E^{(k)}=E_h^{\otimes k},\qquad A^{(k)}=\sum_{j=1}^kA_{h,j},\qquad
 \mathcal I_{h,k,M}=\mathcal P_M^{(k)}
                     \cap(h(s_1),\ldots,h(s_k)),\\
 R_M=(I-P_{\mathcal B_M})s_h^{\otimes k},\quad
 G_M=(J_{k,M}M_{k,M}^{-1}J_{k,M}^*)^{-1},\\
 W_M=-(Y_M^*C_M+C_M^*Y_M),\quad
 G_M-G_{M+1}=Y_M^*Y_M,\quad
 Y_M,C_M:E^{(k)}\longrightarrow\mathcal E_M,\\
 \dim\mathcal E_M=\binom{M+k}{k-1}.
 \end{gathered}
 \tag{R5}
\]
Section~\ref{sec:coherent-tensor-integration} proves the original
ordered division, both cochain signs, every layer map, and the
extra orthogonal projection required when tensor packets are
enlarged.  Its formula (CT.18) calculates the resulting change in
both \(G\) and \(W\).

For a reflection-stable packet, let \(p_M\) be either signed
inertia of \(W_M\), and let \(\mathfrak d_M^{(k)}\) be its
complete tensor departure from (CT.20).  The same source layers give
\[
 \begin{gathered}
 \mathcal D_k\le p_M\epsilon_M,\qquad
 p_M\epsilon_M^2\ge
  2kd^{k-1}\sum_{\rho\in Z}m_\rho
        (\operatorname{Re}\rho-1/2)^2+\mathfrak d_M^{(k)},\\
 p_M\le\min\{\operatorname{rank}Y_M,
              \operatorname{rank}C_M,\lfloor d^k/2\rfloor\},\qquad
 \mathcal D_k^2\le
 \operatorname{Tr}(G_M^{-1}(G_M-G_{M+1}))
 \operatorname{Tr}(G_M^{-1}C_M^*C_M).
 \end{gathered}
 \tag{R6}
\]
The zero-inertia case is proved directly in (CT.25).  The full
ordered sum \(\mathcal D_k\) and all multiplicities are retained
in (CT.20).  The joint density and exact graph closure in
Section~\ref{sec:joint-density-dissipation} show
\(G_M\downarrow0\) for this actual total-degree sequence and
identify every retained jet in the vertical fibre of (JD.11).
The same section proves the coercive inverse-interpolation limit
with its arithmetic unit, the supported-zero limit morphism, and
the accumulated determinant cost (JD.14)--(JD.15).  These formulas
connect the relative layer energy to the exact rate of determinant
loss through the same original matrices.

The September 11 support, Rees, comparison and audit packages
are retained in the companion source directory. The current
analytic baseline is the September 12 sequence: finite operator
comparison (PR~8), chain descent (PR~9), global retraction
(PR~10), homotopy control (PR~11), formal tau recovery (PR~12)
orthogonal boundary control (PR~13), and coherent interpolation
with joint tensor boundaries (PR~14).
Part~II includes the complete selected source proof notes with
their exact revisions. Historical statements that a map remains
unconstructed belong to their stated source dates; a later
explicit construction in this reader or the source chain
supersedes that particular statement.

The new local proofs are the full support-index equivalence,
the corrected carrier and sector comparisons, the arithmetic
confluent-kernel minimum and complete two-row source observation,
the classification of source-metric images and aggregate control,
and the graph/weighted-jet comparison.
The next tranche adds the certified original theta-moment series
and its full truncation bound, the confluent quadrature and
finite/terminal/tail resolvent identities, the explicit initial
convergence rate, and the metric-departure and nilpotent-chain
identities. The recovered June 24 historical packet calculation
is reproduced with a full-jet comparison in
Section~\ref{sec:historical-packet-ghost}. Its exact
\(A\)-invariant saturation is every off-axis local factor;
the quotient, involutions, source-metric transport and split-zero
endomorphism descent are proved there. The original transcript
excerpt and its response locators are retained in the source
package. This is an explicitly recorded historical recovery,
with its new arithmetic saturation proof distinguished from
the original permutation calculation.
The subsequent historical reconstruction in
Section~\ref{historical-hurwitz-reconstruction-and-the-complete-laurent-jet-bridge}
proves the inverse map from a complete local Hurwitz zero-cluster
germ to every moment of its original compact probability measure,
followed by the explicit Bernstein recovery of that measure.
The monic cluster polynomial retains arbitrary zero multiplicity.
The same section proves the complete Laurent period maps and
their collision kernels.  For two distinct primes, the joint
Laurent map is surjective onto the full finite spectral algebra;
Theorem~\ref{thm:two-prime-laurent-presentation} gives every
primary kernel and the inverse original spectral coordinate,
including its truncated logarithm on each nilpotent factor.
The factorwise Split-Zero lifts and their exact comparison with
the synchronized global lift are proved in (S4)--(S5).
Their elementary algebraic and analytic steps are written here.
No claim of global mathematical priority follows from finding
or proving an identity in this research programme.

The finite validation records distinguish exact symbolic
identities, rational calibration models, numerical quadratures
and intentionally rejected negative controls. They do not count
test methods as theorems or imply analytic certification from
an \(L^2\) numerical experiment. The supplied formalization records
are also kept separate: successful exact-head CI metadata
does not turn the unformalized analytic theta arguments into
Lean proofs. No Lean run was performed for this continuation.

The current research calculation includes (R2) on the actual
\(|g/h|^2\)-moments, with its determinant update and two-row
recurrence, together with the topology constants of
Theorem~\ref{thm:weighted-jet}. No uniform arithmetic estimate
forcing (R2) to vanish has been established in this reader.
RH, GRH and a global purity theorem are not conclusions of it.
The exact maps and source constraints above specify the live
route on which tandem work can now continue.
The next estimate must propagate (GC.8) through (KR.16), retaining
\(q_n^2\), \(\kappa_n\), \(a_n\), the terminal ratio \(F_n\)
and the actual tail denominator. The arithmetic moment formulas
give an executable source for those finite data. The exact
departure identity and the full-jet metric bounds specify what
a resulting relative estimate would control in the original
packet; no component is removed by changing observation.

The next step of that propagation is now calculated in
Section~\ref{sec:terminal-continuation}.  For
\(u_n=q_n(A)c\), \(w_n=q_{n-1}(A)c\), and
\(M_n=\kappa_{n-1}^{-1}K_{n-1}^{-1}\), let
\(\beta_n=u_n^*M_nu_n\), and let \(P_n\) be the exact
\(M_n\)-orthogonal projection with kernel \(\mathbb Cu_n\)
when \(\beta_n>0\).  Put
\(p_n=P_nw_n\), \(z_n=P_nAu_n\).  The universal step is
\[
 \begin{gathered}
 M_{n+1}=a_n^{-1}\left(M_n-
      \frac{M_nu_nu_n^*M_n}{a_n+\beta_n}\right),\qquad
 \frac{\det G_{m+1}}{\det G_m}
             =\frac{a_n}{a_n+\beta_n},\\
 \Delta_{n+1}=
 \frac{\beta_n\|z_n+a_np_n\|_n^2}{a_n(a_n+\beta_n)},\qquad
 \epsilon^{(n)}=|g_Z|+\sqrt{g_Z^2+\Delta_n},\quad
 g_Z=\operatorname{Re}\operatorname{tr}A-d/2.
 \end{gathered}
 \tag{R7}
\]
Here \(m=n-d-1\).  The zero-vector exceptional step is given
completely in (KT.14), so critical centres and every polynomial
coincidence remain in (R7)'s source calculation.  The alternate
terminal Gram follows by the exact isometry (KT.15) and recurrence
(KT.16).  The full scalar-to-matrix extraction uses the doubled
off-line algebra and the explicit sequence (KT.23); at a reflected
pair its matrix coefficient is exactly
\((-1)^lF_n^{[k+l+1]}(\rho)\).  Finally (KT.24)--(KT.25)
recover these coefficients from the original initial error and
tail, with both reciprocal series and all polynomial/norm factors
displayed.  The remaining energy estimate is therefore attached to
the three explicit terms of (KT.11), including their actual cross
pairing, on the same arithmetic moment sequence.

For a packet closed under both conjugation and reflection, the
original real-coefficient calculation in (KT.28) closes this step
in terms of consecutive determinant ratios.  With the fixed
section metric \(G_{-1}\) retained, put
\(r_m=\det G_m/\det G_{m-1}\) for \(m\ge0\).
Then the actual arithmetic allowance obeys exactly
\[
 \boxed{\epsilon_m^2
   =a_{d+m+1}\frac{(1-r_m)(1-r_{m+1})}{r_{m+1}}.}
 \tag{R8}
\]
The coefficient reality and full reflection trace prove
\(\gamma_n=0\) in the original coordinates; (KT.7)--(KT.8)
then prove (R8), including zero-step cases.  This places the
relative quartet control directly on the same monic norm ratios
and consecutive Gram determinants as the certified moment problem.

This also gives a direct arithmetic bound for the original tesserine
violation channel.  Retain the original quartet element
\[
 Q(\rho)=\rho\,1+(1-\rho)\sigma
       +\overline\rho\,\kappa+(1-\overline\rho)\sigma\kappa,
 \qquad \rho=\beta+i\gamma,
\]
and its character \(c_{-+}(\rho)=2(2\beta-1)\), as calculated
in the carrier section.  If \(\gamma\ne0\), \(\beta\ne1/2\),
and \(Z\) is this entire four-point zero orbit with common full
order \(\mu\), then \(d=4\mu\) and for every \(m\ge0\),
\[
 \boxed{
 \mu^2|c_{-+}(\rho)|^2
   \le \epsilon_m^2
   =a_{4\mu+m+1}\frac{(1-r_m)(1-r_{m+1})}{r_{m+1}}.}
 \tag{R9}
\]
To prove every coefficient, conjugation and the functional equation
preserve the full order, and the four real displacements are
\(\beta-1/2,\beta-1/2,-(\beta-1/2),-(\beta-1/2)\).
The proved aggregate bound (R3) is therefore
\(\epsilon_m\ge4\mu|\beta-1/2|
=\mu|2(2\beta-1)|\).  Squaring and using (R8) gives (R9).
The sector morphism on the original supported quartet is precisely
\[
 G(\chi_{-+}):G(\mathbb C[K_4])\longrightarrow G(\mathbb C),
 \quad Q(\rho)^\bullet\longmapsto
                   c_{-+}(\rho)^\bullet,\quad\tau\longmapsto\tau.
\]
Here \(\chi_{-+}=\operatorname{ev}_{-1,+1}\) is the original group character extended
linearly; it preserves multiplication because the character does,
and hence its displayed \(G\)-lift is a semiring morphism.
At \(\beta=1/2\) this morphism takes the supported value \(e\);
the same original quartet has phase channel \(c_{--}=4i\gamma\).
The orbit then has two distinct points when \(\gamma\ne0\), so
its degree is \(2\mu\) and (R8) applies with that exact degree;
the left-hand side of the sector bound is zero.  Thus the arithmetic
volume comparison and the supported critical value use the same
typed character map, with the phase coordinate retained.

\subsection{The original theta norm and the positive arithmetic frontier}

The two newly delivered source notes, reproduced completely in the
appendices, admit the following exact common formula with (R8).
Keep \(u_j=(I-P_{L_{j-1}})D^jf_0\),
\(\mathfrak h_j=\|u_j\|^2\), and the original measure
\(|g|^2dt/(2\pi)\).  The polynomial division and original
minimum in (KL.23)--(KL.27) prove, at every full packet,
\[
 \kappa_{d+j}=\mathfrak h_jr_j\quad(j\ge0),\qquad
 \epsilon_m^2=
 \frac{\mathfrak h_{m+1}}{\mathfrak h_m}
       (r_m^{-1}-1)(1-r_{m+1})
       -\mathfrak h_{m+1}^2|c_m|^2
 \tag{R10}
\]
for reflection-stable packets and \(m\ge0\), where
\(c_m=d_mG_m^{-1}d_{m+1}^*\) uses the original theta
projection rows.  The first identity includes \(j=0\) and
the fixed section \(G_{-1}\): the actual new-layer vector
at that endpoint is exactly \(f_0\).  The row dictionary
(KL.30) retains its separate domain \(m\ge0\); its
cross term is \(\mathfrak h_{m+1}c_m=\overline{\gamma_n}\),
with \(n=d+m+1\).  Both signs and the conjugation are
obtained by multiplying those rows against the original inverse
metric, including both local-unit matrices.  Equation (KL.25)
substitutes the first identity in (R10) into the complete terminal
energy, proving its second identity.  Under conjugation closure
the original coefficients are real, so \(\gamma_n\) is real;
the reflection trace makes its real part zero.  Thus \(c_m=0\)
in this case, by precisely the same calculation, with no discarded
cross term.

For the full quartet in (R9), this gives the proved chain
\[
 \boxed{
 \mu^2|c_{-+}(\rho)|^2
 \le\epsilon_m^2
 =\frac{\mathfrak h_{m+1}}{\mathfrak h_m}
              (r_m^{-1}-1)(1-r_{m+1}).}
 \tag{R11}
\]
Indeed the order-preserving symmetries establish the hypotheses of
(R10) and its vanishing cross term; the lower bound was proved
with every multiplicity in (R9).  The norms \(\mathfrak h_j\)
are independent of the packet.  For \(h\mid H\), the
original polynomial map \(P\mapsto(H/h)P\) is an isometry
because \((g/H)(H/h)P=(g/h)P\).  Its jets are the old
complete jets and zero jets at each added centre.  Its fixed
representative and all its corrected representatives agree under
this insertion, as proved in (KL.28).  Thus the common theta
sequence in (R11) connects the actual nested-packet maps; the
packet dependence remains in the two displayed determinant ratios.

For the tensor frontier, retain the original full column maps
\(F,E_+:\mathbb C^r\to E_h^{\otimes k}\) and its original
norm diagonal \(\Omega\) in (AP.4).  With \(G_M\) the
actual total-degree minimum, define
\[
 \begin{gathered}
 \mathsf A=\Omega^{-1/2}F^*G_MF\Omega^{-1/2},\qquad
 \mathsf B=\Omega^{-1/2}E_+^*G_ME_+\Omega^{-1/2},\\
 \Gamma_{k,M}=\max_{|\alpha|=M}
       \sum_{j=1}^k\bigl(a_{\alpha_j+1}+(k-1)a_{\alpha_j}\bigr),
 \qquad \pi_M=\frac{\det G_{M+1}}{\det G_M}.
 \end{gathered}
 \tag{R12}
\]
All complete conjugation-and-reflection packets satisfy
\[
 \boxed{\epsilon_{k,M}^2
   =\lambda_{\max}(\mathsf A^{1/2}\mathsf B\mathsf A^{1/2})
   \le\Gamma_{k,M}\lambda_{\max}(\mathsf A)
   \le\Gamma_{k,M}(\pi_M^{-1}-1).}
 \tag{R13}
\]
Here the square roots are the explicit metric isometries, with
their inverses in (AP.12); the original matrices in (R12) remain
part of the formula.  The full proof in (AP.7)--(AP.21) starts
from the actual linear reflection of the original Mellin source.
It places the two frontiers in opposite orthogonal eigenspaces
and gives \(U^*V=0\).  The relative weight is consequently
the two off-diagonal blocks \(UV^*\) and its adjoint.  Its
complete singular-vector decomposition identifies its norm squared
with the positive operator in (R13), including all zero vectors.
The explicit incidence row sum gives \(VV^*\preceq\Gamma I\);
the inverse-metric update gives
\(\pi_M^{-1}=\det(I+UU^*)\).  Multiplication of its
nonnegative eigenvalue factors proves the two inequalities.
For \(k=1,M=d+m\), (R13), (R11), and (AP.23) calculate
the same original tesserine bound through its norm, volume, and
positive-operator presentations.

\subsection{The symmetric summand, the actual relation layer, and support}

Every factorial in the symmetric comparison is retained.  The
literal orbit-sum injection \(I\), its diagonal
\(D_o=I^*I\), and \(L=D_o^{-1}I^*\) in (AP.24)
give \(G_M^s=I^*G_MI\) and the exact identity
\((G_M^s)^{-1}=LG_M^{-1}L^*\).  The proof uses
permutation invariance of the actual minimum, so that its metric
commutes with the projection \(IL\).  Hence the original
redundant frontier has symmetric maps \(F_s=LF,E_s=LE_+\)
with the unchanged \(\Omega\), and
\[
 \begin{gathered}
 \mathsf A_s=\Omega^{-1/2}F_s^*G_M^sF_s\Omega^{-1/2},
 \qquad
 \mathsf B_s=\Omega^{-1/2}E_s^*G_M^sE_s\Omega^{-1/2},\\
 (\epsilon_M^s)^2
  =\lambda_{\max}(\mathsf A_s^{1/2}\mathsf B_s\mathsf A_s^{1/2})
  \le\Gamma_{k,M}\bigl((\pi_M^s)^{-1}-1\bigr),
 \qquad \pi_M^s=\frac{\det G_{M+1}^s}{\det G_M^s}.
 \end{gathered}
 \tag{R14}
\]
Equations (AP.25)--(AP.36) prove all these maps, the full
positive and zero spectra, the opposite reflection parities,
and the orbit-factorial grouping.  In particular the isometry
\(Q_M=G_M^{1/2}I(G_M^s)^{-1/2}\) satisfies
\(Q_M^*Q_M=I\) and \(Q_MQ_M^*=IL\); it carries
the symmetric relative operator to the invariant restriction of
the original one, retaining the complementary operator on
\(\ker L\).  The continuous cochain projector in
(SS.1)--(SS.9) has exactly this cohomology image.  Its finite
packet version uses the original section and full jet, and all
Koszul signs in that construction remain in its primitive.

The actual next relation layer underlying these matrices is proved
in (KL.7)--(KL.20).  Its complete graph frame, Gram, fixed-order
division primitive, derivative image, and derivative cokernel are
given there.  The comparison with a homogeneous shell is the
exact polynomial identity (KPC.7), with every lower coefficient
in (KPC.6).  Its full graph is the fibre product (KPC.12),
and its two projection kernels are specified in (KPC.11).
For two tensor coordinates the original ratio algebra embeds by
(KPC.16), with exponent
\(\max_{\rho/\sigma=r}(m_\rho+m_\sigma-1)\) at each
coincident ratio.  Its adjacent-shell generator action is
(KPC.22a); the missing second-coordinate multiplier remains in
that formula.  The finite example (KPC.23)--(KPC.28a) calculates
both original shell ranks, all nine node values, and the actual
minimum, establishing the exact scope of that comparison.

For every linear map in these graph, kernel, quotient and
cochain comparisons, the fixed-support lift at the end of the
projective-comparison section maps \((\lambda,v)\) to
\((\lambda,Tv)\) and maps external absence to absence.
Its scalar and additive identities are proved there, including
supported zero.  For measured support, the historical branch
chapter proves the separate full connection from ordinary branches
to frame points through \(\operatorname{Idl}(L)\to L\),
with its exact kernel congruence and descent criterion.  Its
\(L^\infty\) correspondence identifies every Boolean branch
with a multiplicative state and its normal branches with actual
measure atoms.  The factor \(\mu(a)^{-1}\int_a f\,d\mu\)
is retained.  These explicit maps specify the support observations
available to the analytic construction and preserve the nonatomic
support component alongside its atomic quotient.

\subsection{The original derivative cost and both positive-half-line kernels}

The same original derivative layer has the stronger arithmetic identity
proved in (QD.6)--(QD.15). In its exact metric coordinates,
\(-\widehat Y^*\widehat C=UV^*\) occupies one off-diagonal
reflection block. Thus its norm is the full relative allowance,
and the original trace cost \(\chi_M\), operator cost \(\zeta_M\)
and largest update eigenvalue \(\lambda_M\) satisfy
\[
 \epsilon_{k,M}^2\leq
        \frac{\lambda_M}{1+\lambda_M}\zeta_M
 \leq\frac{\lambda_M}{1+\lambda_M}\chi_M,
 \qquad
 \epsilon_{1,M}^2=(1-\pi_M)\chi_M.
 \tag{R15}
\]
The first two inequalities follow from the norm of the product of
the two actual layer maps and the eigenvalues of
\(UU^*(I+UU^*)^{-1}\). For one factor, the layer is
one-dimensional, so that product is the outer product of its two
original rows and its norm is exactly their norm product. Its
rank-one determinant update is \(\pi_M=(1+\lambda_M)^{-1}\),
proving the last equality. This argument includes either row zero.
The full lower bounds \(4k^2\delta_*^2\) and all finite accumulated
volume constraints follow in (QD.10)--(QD.13), with their positive
denominators established before division.

For one factor, put \(z=s-1/2\), \(x_h=-z^2\), and retain the
monic polynomial \(H\) defined by \(h(s)=(-1)^DH(x_h)\),
\(d=2D\). The exact algebra is
\(A_H\oplus zA_H=\mathbb C[x_h,z]/(H,z^2+x_h)\), with its
two local length-\(m_\rho\) branches explicitly inverted in
(SP.8)--(SP.9). The original source is the Hilbert direct sum for
\(d\lambda_h=|V(x_h)|^2dx_h/(2\pi\sqrt{x_h})\) and
\(x_h d\lambda_h\). Here \(v_h(s)=V(x_h)\) is the entire
original quotient and its full unit modulo \(H\) enters both
kernels. With those kernels and their inverse metrics, (SP.23)--(SP.35)
give
\[
 \begin{gathered}
 \epsilon_{N-d}=
  \left\|G_e^{-1/2}(G_o-G_eX)G_o^{-1/2}\right\|,\\
 K_e-XK_o=
 \begin{cases}f_ne_n^*/p_n,&N=2n,\\
              -f_{n+1}e_n^*/u_n,&N=2n+1.
 \end{cases}
 \end{gathered}
 \tag{R16}
\]
The first line follows by substituting the exact multiplication
matrix \(\frac12I+\left(\begin{smallmatrix}0&-X\\I&0\end{smallmatrix}\right)\)
in the original weight form. The second is its full inverse-metric
recurrence: all interior adjacent terms cancel with their original
norm ratios, and the last pair has the two displayed parity signs.
The finite kernels have indices \((n,n-1)\) and \((n,n)\),
respectively. Hence the norm is given by the full quadratic
contractions in (SP.33), with at most one nonzero singular value.
Every nilpotent jet remains inside \(f_n,e_n,X\) and both metrics.
The exact derivative cost (SP.34) identifies (R16) with (R15),
including the two zero-frontier cases.

\subsection{The gamma determinant and the retained spectral-sum fibre}

The original completed function has the exact factorization
\(g(1/2+it)=\Gamma(1/4+it/2)a(t)\), with
\(a(t)=-\pi^{-1/4}(t^2+1/4)e^{-it\log\pi/2}\zeta(1/2+it)\).
The beta integral and generating-function calculation (TG.7)--(TG.14)
prove the reference monic norms
\(\gamma_n=\sqrt2\,n!\Gamma(n+1/2)\). Keeping the actual
multiplier in the full positive arithmetic Gram matrix gives its
determinant \(D_n\), with \(D_{-1}=1\). Monic triangular Gram
congruence proves \(\mathfrak h_n=\gamma_nD_n/D_{n-1}\), so
the actual quartet formula is exactly
\[
 \boxed{\epsilon_m^2=(m+1)(m+\tfrac12)
        \frac{D_{m+1}D_{m-1}}{D_m^2}
                  (r_m^{-1}-1)(1-r_{m+1}),\qquad m\ge0.}
 \tag{R17}
\]
This follows by substituting the consecutive norm ratio into
(R10), whose cross term vanishes by the proved quartet symmetry.
The full packet determinants are independently expressed in
(TG.17)--(TG.18), and the original confluent kernel is the complete
Gram congruence (TG.20)--(TG.22). The explicit triangular maps to
the parameter \(3/4\) family and their inverses retain every phase
and degree. The theorem-hypothesis analysis (TG.25)--(TG.28)
proves the obstruction from the infinitely many retained critical
zeros on positive-measure intervals; it also supplies these exact
Gram maps for continuing beyond that obstruction.

The web spectral-sum construction retains an additional relative
coordinate \(\Delta=(s_1-s_2)^2\). Write \(S=s_1+s_2\),
\(Z=S-1\) and \(x_\Sigma=-Z^2\). Its complete reflection
components, proved in Section~\ref{sec:spectral-sum-stieltjes-pair}, are
\[
 B_e=\frac{\mathbb C[x_\Sigma,\Delta]}{(A_0,x_\Sigma\Delta A_1)},
 \qquad
 B_o=\frac{\mathbb C[x_\Sigma,\Delta]}{(A_0,\Delta A_1)}.
 \tag{R18}
\]
Here the original polynomial identities are
\(H_0(S,\Delta)=A_0(x_\Sigma,\Delta)\) and
\(H_1(S,\Delta)=Z A_1(x_\Sigma,\Delta)\); all coefficients
come from \(h((S+r)/2)=H_0+rH_1\), \(r^2=\Delta\).
The map \(\alpha:B_o\to B_e\) multiplies by \(x_\Sigma\),
and \(\beta:B_e\to B_o\) is the quotient. They satisfy
\(\alpha\beta=x_\Sigma\) and \(\beta\alpha=x_\Sigma\).
The original generator is
\(S=I+\left(\begin{smallmatrix}0&-\alpha\\\beta&0\end{smallmatrix}\right)\).
Their dimensions are \(D(D+1)\) and \(D^2\); both values follow
by taking the symmetric square of the original two
\(D\)-dimensional reflection eigenspaces. The exact maps, rather
than an identification of these unequal spaces, enter the weight.

The full kernel and cokernel of \(S-1\) both recover \(A_H\).
Indeed \(A_0(0,\Delta)=(-1)^DH(-\Delta/4)\ne0\) makes
\(\alpha\) injective, while \(\beta\) is onto. Its cokernel
map evaluates \(x_\Sigma=0\) and sends \(\Delta\) to
\(-4x_h\). The kernel injection sends
\(q(x_h)\) to \(\Delta A_1\,q(-\Delta/4)\) in \(B_e\),
with the full polynomial \(\Delta A_1\) retained before taking
the quotient. The ideal calculation and the dimension count prove
both exact sequences in the cited section, including all orders.
On both defect modules the actual tensor unit acts as
\(V(x_h)^2\), because at \(S=1\) the two original arguments
are \(s\) and \(1-s\). Thus the complete paired packet, its
nilpotents and its arithmetic unit remain in the relative fibre.

The exact matrix-weight pushforward retains all powers of
\(\Delta\), the Jacobian \(dx_\Sigma/\sqrt{x_\Sigma}\),
and both vector Hilbert measures. Its actual least-norm quotient
gives the two metrics in (R18); the upper block of the original
weight is \(\beta^*G_o-G_e\alpha\), with its full rectangular
singular spectrum. All maps preserve external absence and carry a
vanishing amplitude to its supported zero through the displayed
fixed-support lifts. These finite identities and complete source
maps are established without an asymptotic estimate for (R17).

\subsection{The sum derivative, its strict normal energy and the full conormal algebra}

For the actual sum amplitude retain
\(a_h^{\rm SC}(t)=v_h(1/2+it)/\sqrt{2\pi}\),
\(w=|a_h^{\rm SC}|^2\), \(\mu_h=\int w\), and
\(I_h=4\int|(a_h^{\rm SC})'|^2\).
In the original coordinates \(t_i=u/k+y_i\),
\(y_k=-\sum_{i<k}y_i\), the full product amplitude is
\(\Psi=\prod_i a_h^{\rm SC}(t_i)\) and its squared fibre norm
is \(m=w^{*k}\). The exact Jacobian has absolute value one.
Differentiating before integration gives
\(\partial_u\Psi=k^{-1}\sum_i a_i'\prod_{j\ne i}a_j\).
The mixed integrated terms vanish because
\(\int\overline a\,a'=0\). Orthogonal projection onto the
actual product-amplitude line gives
\[
 \int_{\mathbb R}\frac{|m'|^2}{m}\,du
       +4\int_{\mathbb R}\|n(u)\|^2\,du
          =\frac{\mu_h^{k-1}}k I_h,\qquad
 n=\partial_u\Psi-\frac{m'}{2m}\Psi .
 \tag{R19}
\]
The proof in (FC.1)--(FC.4) retains the original mass and
phase, and establishes every required integrability statement.
For every finite independent relative polynomial frame,
\(j_uc=\Psi\sum_\alpha\theta_\alpha(i\mathbf y)c_\alpha\),
the corresponding exact derivative is
\[
 \partial_u(j_uc)=j_u(c'+\Gamma_uc)+N_uc,\qquad
 \Gamma=W^{-1}j^*j',\quad
 N=(I-jW^{-1}j^*)j',\quad j^*N=0 .
 \tag{R20}
\]
This identity follows by inserting the orthogonal projection onto
\(\operatorname{im}j_u\), with \(W=j^*j>0\). Its full normal
Gram is \(\mathcal N=N^*N\).

Every finite packet leaves infinitely many critical-line zeros of
the original amplitude. At such a zero \(T\) of complete order
\(\ell_T\), the transverse expansion in (FC.25) has coefficients
\(c_TA_b\delta^{\ell_T}\) for \(\Psi\) and
\((\ell_T/k)c_TA_b\delta^{\ell_T-1}\) for \(\partial_u\Psi\),
with \(c_TA_b\ne0\). Thus an equality \(N_uc=0\) forces the
relative polynomial to vanish on the entire hyperplane
\(y_1=T-u/k\) whenever the other factors can be chosen nonzero.
For \(k\ge3\) this holds at every \(u\); for \(k=2\) it holds
outside the countable sumset \(\mathcal Z_h+\mathcal Z_h\).
Infinitely many such hyperplanes force the finite polynomial to
vanish. Consequently (FC.21)--(FC.27) prove
\[
 \begin{gathered}
 \mathcal N(u)>0\ \text{everywhere for }k\ge3,\\
 \mathcal N(u)>0\ \text{almost everywhere for }k=2,\qquad
 \int\frac{|m'|^2}{m}<\frac{\mu_h^{k-1}}kI_h .
 \end{gathered}
 \tag{R21}
\]
The full proof retains the orders, coefficient domains and the
quartet two-factor zero fibre \(N_0=0\).

The two-branch transformation is equally exact for a matrix
weight which is not even. At \(x>0\), put \(v=\sqrt x\) and
\[
 M=\begin{pmatrix}I&ivI\\I&-ivI\end{pmatrix},\qquad
 H=\frac{M^*\operatorname{diag}(W(v),W(-v))M}{2v}.
\]
The original coefficient is \(c(\pm v)=a(x)\pm ivb(x)\).
Equations (FC.10)--(FC.18) prove the full unitary correspondence,
the covariant derivative \(\nabla_x=\partial_x+\mathcal A\),
and its energy
\[
 \begin{gathered}
 \mathcal A=M^{-1}
   \operatorname{diag}\!\left(\frac{\Gamma(v)}{2v},
                         -\frac{\Gamma(-v)}{2v}\right)M
                    +\operatorname{diag}(0,I/(2x)),\\
 \int\|\partial_u(jc)\|^2\,du
   =4\int_0^\infty x(\nabla_xq)^*H\nabla_xq\,dx
                      +\int_0^\infty q^*\mathsf Qq\,dx ,
 \qquad q=(a,b).
 \end{gathered}
 \tag{R22}
\]
Here \(\mathsf Q=M^*\operatorname{diag}(\mathcal N(v),
\mathcal N(-v))M/(2v)\). The retained Jacobian drift is
\(\mathcal A^*H+H\mathcal A=H'+H/(2x)\).
The two endpoint traces \(a+i\sqrt x\,b\) and
\(a-i\sqrt x\,b\) must agree at \(x=0\); this is the proved
joining condition for the original derivative domain.

The arithmetic derivative has the unital algebra bridge
\[
 \Phi:P/I^2\longrightarrow (P/I)[\epsilon]/(\epsilon^2),
 \qquad [p]\longmapsto[p]+\epsilon[\partial_\Sigma p],
 \quad I=(h(s_1),\ldots,h(s_k)),\quad
 \partial_\Sigma=k^{-1}\sum_i\partial_{s_i}.
 \tag{R23}
\]
The product rule modulo \(I\) proves multiplicativity. Its
conormal restriction is the row
\((a_i)\mapsto k^{-1}\sum_i h_i'a_i\) on
\(I/I^2\cong(P/I)^k\). On an ordered local block of lengths
\(m_i\), multiplication by \(h_i'\) sends the input constant
in variable \(i\) to
\(m_ic_i(0)z_i^{m_i-1}\) times that constant. Counting the
monomials having at least one top exponent proves its full rank
\(d^k-(d-r_0)^k\), \(r_0\) the number of distinct packet roots.
The kernel and image of the algebra map, including its
linear cokernel, are proved in (CW.6)--(CW.11).

For the original monic remainder section \(j:P/I\to P/I^2\),
\(Z=S-k/2\) and \(\mathsf D_0=\delta j\), the exact section
defect and correction are
\[
 [Z,\mathsf D_0]=-I+\mathsf J,\qquad
 \mathsf J=\frac1k\sum_i M_{h_i'}c_i\ell_i
          =\frac d k\sum_iP_i,\qquad
 P_i=\frac1dM_{h_i'}c_i\ell_i .
 \tag{R24}
\]
Here \(\ell_i\) extracts the original \(s_i^{d-1}\) coefficient
and \(c_i\) inserts the constant polynomial in that variable.
The identity \(\ell_iM_{h_i'}c_i=dI\) proves that the \(P_i\)
are commuting idempotents. Their complete tensor decomposition
gives eigenvalues \(dj/k\) with multiplicities
\(\binom kj(d-1)^{k-j}\) for \(d>1\), and
\(\mathsf J=I\) for \(d=1\).
Its trace is \(d^k\) and its rank \(d^k-(d-1)^k\).
The explicit factorization (CW.22) through the full conormal row
proves the relationship between these two rank formulas.

The complete arithmetic unit
\(\widehat U=[\prod_i v_h(s_i)]_{I^2}\) gives
\(j_U=\widehat UjU^{-1}\),
\(\mathsf D_U=M_{U^{-1}\delta\widehat U}
+M_U\mathsf D_0M_U^{-1}\), and
\([Z,\mathsf D_U]=-I+U\mathsf JU^{-1}\).
For the actual minimum section \(j_M\), the difference
\(a_M=j_M-j_U\) is its admitted original relation in \(I/I^2\);
the correction is exactly
\(U\mathsf JU^{-1}+[Z,\delta a_M]\), with trace \(d^k\).
These are the full source maps proved in (CW.23)--(CW.30).
In one variable they also give
\[
 \mathsf J_Uu=(j_hg')\,\mathscr R_{\mathcal Z}(1,u),\qquad
 \mathscr R_{\mathcal Z}(1,u)=\ell(U^{-1}u),\qquad
 \mathscr R_{\mathcal Z}(1,j_hg')=\sum_\rho m_\rho=d .
 \tag{R25}
\]
The partial-fraction proof in (CW.32)--(CW.36) retains the
full poles before taking residues. Finally \(G(\Phi)\) is the
proved semiring lift; the derivative is its infinitesimal
coefficient map. Thus every relation killed by the quotient
retains its supported zero and its explicit derivative observation.

\subsection{The full relative limit and the original arithmetic observation}

Retain $k\ge2$ and the nonempty original packet $d=\deg h\ge1$
throughout this subsection. For two nested relative frames, keep the literal coefficient inclusion
$E_{LM}$ and the original matrices of (NT.1)--(NT.37). Their exact
connection difference $K_{LM}=\Gamma_LE_{LM}-E_{LM}\Gamma_M$ obeys
\[
 N_M=j_LK_{LM}+N_LE_{LM},\qquad
 N_M^*N_M=K_{LM}^*W_LK_{LM}+E_{LM}^*N_L^*N_LE_{LM}.
 \tag{R26}
\]
The cross terms vanish because $j_L^*N_L=0$. Thus enlarging a frame
transfers the first displayed component into its tangent part, with its
complete positive Gram retained. The arithmetic exponential bounds and
weighted polynomial density proved in (ND.1)--(ND.34) give
$\Pi_M\to I$ strongly on every original relative fibre. For each fixed
original source polynomial $P$, (NN.1)--(NN.22) prove
$r_M(P)=(I-\Pi_M)\partial_u(\Psi P)\to0$ in the original global
$L^2$ norm. Their normal-image transitions are the explicit contractions
$r_M(P)\mapsto(I-\Pi_L)r_M(P)=r_L(P)$; their full derivative-graph
transitions keep $\partial_u(\Psi P)$ unchanged and are isometric.

The same proofs retain the actual logarithmic arithmetic observation.
The invariant relation $P_{\rm rel}=\sum_i h(s_i)$ has zero ordinary
jet but derivative jet $(U_k/k)\sum_i[h'(s_i)]\ne0$.
For $d=1$ the sum is $k$; for $d>1$ its distinct highest standard
monomials have coefficients $d$, so it cannot vanish. Its normal
seminorm tends to zero, while its algebraic observation remains this
specified nonzero vector. The exact split quotient of that zero limiting
seminorm sends every supported source to $e$ and external absence to
$\tau$; the nonzero observation therefore cannot factor through that
quotient. Its original finite maps and full derivative graph remain.
For the changing canonical target metrics the complete comparison is
\[
 G_M\preceq G_*,\quad
 B_M=G_*^{-1/2}G_MG_*^{-1/2},\quad
 x^*G_Mx=(G_*^{1/2}x)^*B_M(G_*^{1/2}x).
 \tag{R27}
\]
The fixed-source limit has not estimated this changing $B_M$.
For $h=1$ the finite coefficient quotient is zero, so the arithmetic
observation above is its zero map; the original analytic source and
the fixed-source density and derivative-graph maps remain as in (NN.22).
The cyclic/full-frame derivative split and canonical-section difference
are proved with all four metric terms in (CM.1)--(CM.19).

\subsection{Every relation depth and its full cyclic complement}

For the same monic $h$ and $I=(h(s_1),\ldots,h(s_k))$, the complete
depth calculation (CC.1)--(CC.52) proves
\[
 \begin{gathered}
 I^r\cap\mathbb C[S]=(\chi_r),\qquad
 \chi_r(T)=\prod_\lambda(T-\lambda)^{\ell_{\lambda,r}},\\
 \ell_{\lambda,r}=\max_{\sum_i\rho_i=\lambda}
 \left(1+\sum_i(m_{\rho_i}-1)+(r-1)\max_i m_{\rho_i}\right),\\
 C_r=\mathbb C[T]/\chi_r\xrightarrow{\eta_r=U_r\alpha_r}
 A_r=(P/I^r)^{\mathfrak S_k},\quad
 \Pi_r\eta_r=I,\quad A_r\cong C_r\oplus\ker\Pi_r.
\end{gathered}
 \tag{R28}
\]
The coefficient-extraction formula (CC.13) constructs this retraction,
and (CC.10) counts every invariant orbit before taking its trace.
The quotient derivative $\delta_r:P/I^{r+1}\to P/I^r$ and its scalar
counterpart satisfy the commuting squares (CC.17). Their full unit
term has the proved decomposition
$\delta_r\eta_{r+1}=\eta_r(\partial_r+M_{b_r}\rho_r)
+\mathcal N_r\rho_r$ of (CC.23); its complementary quotient map
vanishes exactly when $U_r^{-1}\delta_rU_{r+1}\in\alpha_rC_r$.
The first-jet maps, their full kernels, the product-Jacobian
factorization, and the exact two-factor parity ideals at every depth
are proved in (CC.29)--(CC.51b). An arbitrary weighted multiplier
has all four blocks (CC.45); its lower diagonal block is a genuine
quotient operator precisely when the lower left block is zero.
Thus the earlier weighted-trace language has the exact operator
interpretation (CC.45)--(CC.49), with no complementary direction lost.

\subsection{Exterior amplification and the arithmetic volume to estimate}

For the actual cyclic quotient, keep $A=M_S$, its canonical $G_N$ and
$H=G_N^{-1}(A^*G_N+G_NA-kG_N)$. The complete exterior source note
proves that $H$ has spectrum $\epsilon,-\epsilon,0,\ldots,0$.
The additive exterior action $A^{[p]}$ and the original unscaled
representative $R_N^{\otimes p}\operatorname{Alt}_p$ have source Gram
$p!\bigl(\det(G_N)_{I,J}\bigr)_{I,J}$; their control has the same
literal factorial and allowance $\epsilon$. On the determinant of
the full positive generalized spectral subspace this gives
\[
 L_{h,k}=\sum_{\Re\lambda>k/2}\ell_\lambda(2\Re\lambda-k)
 \le\epsilon_{h,k,N}^{\rm cyc}.
 \tag{R29}
\]
The full cochain primitive carries the sign $(-1)^{k(j-1)}$ in
block $j$, which the tensor differential repeats; their product is
one. Hence this exterior comparison uses the original theta boundary.
The cochain idempotent splits full cohomology, while the ordinary
trace in (R29) is on the specified finite arithmetic packet.

For the exact complete nonreal quartet polynomial of (E33), with
common order $m$, let $\ell_k=1+k(m-1)$. The counted sum grid yields
\[
 q=\ell_k(k+1)^2,\quad
 L_{h,k}=2\delta\ell_k(k+1)
 \left\lfloor\frac{(k+1)^2}{4}\right\rfloor,\quad N\ge q-1.
 \tag{R30}
\]
The explicit four occupation counts in (E33)--(E34) prove this formula
without a noncollision assumption on ordered tuples. In a larger
packet it remains a lower bound and the threshold is its actual
$\deg\chi_{h,k}-1$. The equality analysis (E1)--(E36) further proves,
in the original metric, with positive spectral projector $Q$,
orthogonal range projection $P$, and $T=PA(1-P)$,
\[
 \begin{gathered}
 \epsilon^2-L^2=\|T\|_{\mathrm{HS},G_N}^2
       +\epsilon^2\{\det M+\det(I_2-M)\},\\
 BX-XD=T,\quad X=(Q-P)|_{\ker P},\quad
 X=\int_0^\infty e^{-tB}T e^{tD}\,dt.
\end{gathered}
 \tag{R31}
\]
Here $M$ is the explicit positive compression in (E5)--(E7),
and $B,D$ are the original diagonal blocks of $A$ in the orthogonal
sum. The finite nilpotent exponential formulas (E15) prove convergence
of the integral and that this Sylvester map is invertible.
For $L>0$, equality $L=\epsilon$ is equivalent to $P=Q$, or to
$T=0$. A retained critical Jordan block forces strict inequality.
The canonical Gaussian equality calibration (E35)--(E36) proves why
positivity of a moment measure alone supplies no further strictness.

Finally the arithmetic matrices themselves give the exact calculation
(CV.1)--(CV.12). Set $D_N=\det K_N$ and
$\widetilde K_N=K_{N-1}+b_{N+1}b_{N+1}^*/\omega_{N+1}$.
Then
\[
 \begin{gathered}
 \epsilon^2=\frac{\omega_{N+1}}{\omega_ND_N}\Delta_N,\qquad
 \Delta_N=D_{N+1}+D_{N-1}-D_N-\det\widetilde K_N,\\
 \Delta_N=\sum_{\substack{I\subset\{0,\ldots,N-1\}\\|I|=q-2}}
 \frac{|\det[b_I,b_N,b_{N+1}]|^2}
 {\omega_N\omega_{N+1}\prod_{i\in I}\omega_i}\quad(q\ge2).
\end{gathered}
 \tag{R32}
\]
Determinant multilinearity proves the positive sum, including the
singular first-degree boundary, and (CV.10a) constructs its original
source vector with norm squared $\Delta_N$ and tensor norm squared
$q!\Delta_N$. At $N=q-1$, (CV.11)--(CV.12) give the same control as
$\sum_{j<q-1}|\gamma_j|^2\omega_j/\omega_{q-1}$ in the unchanged
orthogonal expansion of $\chi$. Thus the current quantitative
expression is $\omega_{N+1}\Delta_N/(k^6\omega_ND_N)$ along
admitted growing degrees. Its limit has not been estimated. The
complete lower bound (R30), the positive exterior data (R32), and
the changing metric (R27) are retained together for that calculation.

\subsection{The original Toda volumes and the complete four-term loss}

The Toda continuation now calculates this same quantity from the original
source and relation volumes. In the notation proved in (TVB.1)--(TVB.38c),
use $k\ge1$, a nonempty dagger-stable packet with its full zero orders,
$q\ge1$, real $|\theta|<\pi/2$, and $N\ge q$, except at the
separately stated $N=q-1$ endpoint, throughout (R33)--(R34). We
retain $c=k/2$, $Z(\theta)=M_h(\theta)^k$ and
$Z_\chi=\overline\chi(c-i\partial_\theta)\chi(c+i\partial_\theta)Z$.
The Hankel determinants $\mathfrak D_n=\det[Z^{(a+b)}]_{a,b<n}$ and
$\mathfrak B_m=\det[Z_\chi^{(a+b)}]_{a,b<m}$ are the actual source
and relation Gram determinants. With $m=N-q+1$ their exact determinant-line
map and the rank-one update give
\[
\begin{gathered}
 V_N=\det G_N=\frac{\mathfrak D_{N+1}}{\mathfrak B_{N-q+1}},
 \qquad \delta_N=\frac{V_N}{V_{N-1}},\qquad
 \alpha_{N+1}=\frac{\omega_{N+1}}{\omega_N},\\
 \psi_N=\frac{\operatorname{Tr}A-cq}{i}-\partial_\theta\log V_N,
 \qquad \frac{b_N^*G_Nb_{N+1}}{\omega_N}=i\psi_N,\\
 \epsilon_N^2=
 \alpha_{N+1}(1-\delta_N)(\delta_{N+1}^{-1}-1)-\psi_N^2
 =\alpha_{N+1}\frac{\Delta_N}{\det K_N},\qquad N\ge q.
\end{gathered}
\tag{R33}
\]
Every equality has its complete source, phase and determinant proof in
the cited bridge. At $N=q-1$, the exact replacement remains
$\epsilon_{q-1}^2=(\nu_0-\omega_q)/\omega_{q-1}-\psi_{q-1}^2$,
with $\nu_0=\|\chi\|^2$; no preceding inverse occurs. The explicit
conjugate-linear map (TVB.38b) sends the lower orthogonal projection of
$\chi$ to the actual positive-minor exterior source vector, with the
full squared-norm factor $\det K_{q-1}/\omega_q$.

The two scalar Toda curvatures have full source maps. For $X=(S-c)/i$,
the actual least-norm lift satisfies $R_N'=-I_N$, where
$I_N=P_{\mathcal D_N}XR_N$ and $O_N=(1-P_N)XR_N$ in the current
tilted pairing, where $P_N$ is the tilted source projection onto
$\mathcal P_N$. With $\beta_m=\nu_m/\nu_{m-1}$ for $m\ge1$, the
proofs of (TVB.17)--(TVB.21) give
\[
\begin{gathered}
 G_N''-G_N'G_N^{-1}G_N'=O_N^*O_N-I_N^*I_N,\\
 \operatorname{Tr}(G_N^{-1}O_N^*O_N)=\alpha_{N+1}(1-\delta_N),\qquad
 \operatorname{Tr}(G_N^{-1}I_N^*I_N)=\beta_m-\alpha_{N+1}\delta_N.
\end{gathered}
\tag{R34}
\]
The common subtracted term is present in both positive costs. At the
first admitted degree the two costs are $\alpha_q$ and zero, with no
negative-index relation norm. Every derivative boundary has the literal
cochain primitive (TVB.9), including its tensor signs.

The coordinating Zeta contribution joins the two volume losses with the
original spectral coupling. At $\theta=0$, $N\ge q$ and $L>0$, let
$\mathsf P$ be the $G_N$-orthogonal projection onto the full positive
spectral subspace, $\mathsf C=\mathsf PA(1-\mathsf P)$ its restricted
coupling, and $M=E^\sharp\mathsf PE$, where $E$ embeds the two
$G_N$-unit eigenvectors of $A^\sharp+A-kI$ at $\pm\epsilon_N$.
Set $v_N=-\log\delta_N$, $s_N=(v_N+v_{N+1})/2$ and
$t_N=(v_{N+1}-v_N)/2$. The complete derivation in
(TVB.39)--(TVB.43), using the same metric on both sides, proves
\[
\begin{aligned}
 \alpha_{N+1}\sinh^2s_N-L^2
 ={}&\alpha_{N+1}(e^{t_N}-\cosh s_N)^2+\psi_N^2\\
 &+\|\mathsf C\|_{\mathrm{HS},G_N}^2
  +\epsilon_N^2\{\det M+\det(I_2-M)\}.
\end{aligned}
\tag{R35}
\]
All four terms are nonnegative. In particular the two-step volume lower
bound retains $L^2+\psi_N^2+\|\mathsf C\|_{\mathrm{HS},G_N}^2$
inside its inverse-hyperbolic-sine expression. The exact Sylvester inverse
maps $\mathsf C$ to the difference between spectral and orthogonal
projectors; its full original nilpotent factors and singular value remain
specified. The complete supplied join is also reproduced as a source appendix.

\subsection{Certified theta input and a calculated comparison of the two gaps}

The theta continuation proves the sector construction and the exact maps
\[
 h(D)F_h=f_0,\qquad
 w_1(t)=|h(1/2+it)|^2w_h(t),\qquad
 M_h(\theta)=2\int_1^\infty|F_h(e^{i\theta/2}x)|^2\,dx.
\tag{R36}
\]
The absolute-square integral is for real $|\theta|<\pi/2$; its
holomorphic strip continuation is proved in the theta input chapter.
The dagger sign, complex Fourier phase, endpoint Jacobian and source
factor $g=2\xi$ are all retained in (TI.1)--(TI.16). The exact
incomplete-gamma double series and its explicit infinite-square-tail
bounds (TI.19)--(TI.30) supply rigorous interval enclosures for
$M_1(0)$, $M_1(1/10)$, $M_1''(0)$ and
$a_1^{(k)}/k=M_1''(0)/M_1(0)$. Each displayed rational interval in
that proof has width $3\cdot10^{-40}$. The certificate adds the proved
infinite tail to outward-rounded ball arithmetic; it does not infer
an enclosure from precision agreement. The source norm is positive
for $h=1$ while its arithmetic cyclic quotient is zero. The density
identity in (R36), and the globally constructed Mellin division, specify
the exact passage to packet input without division of a measured zero
value by zero. No packet determinant estimate follows from these four
seed intervals alone.
For the empty packet $h=1$, the arithmetic quotient is zero and
$V_N=1$ by the empty-determinant convention. For $q=1$,
$\epsilon_N=0$ at every admitted degree, whereas the upper Toda
expression is strictly positive for $N\ge1$, as proved in the bridge.

Finally (TC1)--(TC13) calculate both gaps in the original Gaussian
calibration with mass $\mu$, variance $\sigma^2$, $c=1/2$ and
$\chi=(S-c)^2-\sigma^2$. The source and quotient coordinate maps are
fully displayed there. They give
\[
\begin{gathered}
 (V_1,V_2,V_3)=\mu^2\sigma^2(1,1/3,1/11),\quad
 \epsilon_1^2=L_{\mathrm{spec}}^2=4\sigma^2,\quad
 \epsilon_2^2=16\sigma^2/3,\\
 \alpha_3\frac{(\mathcal R_2-1)^2}{4\mathcal R_2}-\epsilon_2^2
 =49\sigma^2/33,\qquad
 \epsilon_2^2-L_{\mathrm{spec}}^2
 =L_{\mathrm{spec}}^2\|Q-P\|_{\mathrm{HS},G_2}^2=4\sigma^2/3.
\end{gathered}
\tag{R37}
\]
The first gap is precisely the Toda imbalance square; the second has
the explicitly calculated spectral/orthogonal projector map. This
calibration uses declared polynomial coefficients. The actual growing
packet calculation continues with (R33)--(R36), retaining its arithmetic
measure and all four terms of (R35).


\subsection{Gamma convolution, certified higher input, and the complete quotient comparison}

\paragraph{Result R38: an exact analytic and finite coefficient interface.}
Keep the original $w_h(t)=|g(1/2+it)/h(1/2+it)|^2/(2\pi)$ and
$r_\lambda(t)=|\Gamma(\lambda+it/2)|^2/(2\pi)$, with
$c_\lambda=2^{1-2\lambda}\Gamma(2\lambda)$. The full Gamma source
appendix and the coefficient-join proof construct
\[
 r_{\lambda,k}=\frac{c_\lambda^k}{c_{k\lambda}}r_{k\lambda},\qquad
 m_{h,k}=w_h^{*k}=r_{\lambda,k}C_\Sigma(B_{h,\lambda}^{\otimes k}),
 \quad B_{h,\lambda}=|v_h/\Gamma(\lambda+it/2)|^2.
\]
The Hilbert maps $U_\Sigma f=f(\sum t_i)$ and its conditional adjoint
$C_\Sigma$ have the inverse paired decomposition
$F\leftrightarrow(C_\Sigma F,(I-U_\Sigma C_\Sigma)F)$.
All masses and the relative vector remain. With $\alpha=2\lambda$,
\[
 c_{h,j}=\frac{\int b_j^{(\lambda)}w_h}{c_\lambda j!(\alpha)_j},\quad
 \mathcal A_h(z)=\sum_{j\ge0}(\alpha)_jc_{h,j}z^j,\quad
 d_{h,k,n}=[z^n]\mathcal A_h(z)^k,
\]
the exact finite moment is
$\int b_n^{(k\lambda)}m_{h,k}=c_\lambda^kn!d_{h,k,n}$.
The original degree-$N$ source and relation Grams at theta zero use
$c_{h,0},\ldots,c_{h,2N}$; their derivative of order $r$ uses the
list through $2N+r$, through the fully written original-$S$ polynomial
test. The analytic join supplied by the Zeta task and proved in the
coefficient chapter is
\[
 \mathcal A_h(z)=c_\lambda^{-1}(1+z^2)^{-\lambda}
                   M_h(\arctan z),\qquad |z|<1,
 \quad M_h(\theta)=\int e^{\theta t}w_h(t)\,dt.
\]
The specified branches at zero and exponential domination justify
all coefficient derivatives. Its inverse uses
$M_h(\theta)=c_\lambda\sec^{2\lambda}\theta\,
\mathcal A_h(\tan\theta)$ for $|\operatorname{Re}\theta|<\pi/4$,
with the analytic branches continued from zero.
The complete source appendix is preserved alongside the explicit
cochain correction; applying $\mathcal V$ to a section difference
gives the boundary, and (GP.26)--(GP.27) construct its primitive.

\paragraph{Result R39: certified original moments and signed coefficients.}
Equations (GS.1)--(GS.40) prove the original theta derivative recurrence
\begin{gather*}
 P_0(y)=8y^2-12y,\qquad
 P_{r+1}(y)=-2yP_r'(y)+(2y-\tfrac12)P_r(y),\\
 \mu_{2r}=2\int_1^\infty
 \left|\sum_{n\ge1}P_r(\pi n^2x^2)e^{-\pi n^2x^2}\right|^2dx.
\end{gather*}
Their fully specified positive square tail (GS.18) bounds the
omitted infinite integer sum at every fixed order. At order two
and three, the rigorous rational enclosures are
\begin{align*}
 \mu_4\in[&383.2743417173642425544072374791308324300288,\\
          &383.2743417173642425544072374791308324300291],\\
 \mu_6\in[&17946.5273016052076235667646919834204772848931,\\
          &17946.5273016052076235667646919834204772848934].
\end{align*}
The original reference $c_{13/4}=10395\sqrt\pi/(2048\sqrt2)$ is
enclosed separately from the actual mass $\mu_0$.
The complete invertible moment/coefficient map (GS.34)--(GS.35)
certifies $c_0,c_2,c_4,c_6$ and proves $c_4<0,c_6<0$.
Its tensor coefficient of degree six retains all three terms
\[
 d_{k,6}=k c_0^{k-1}a_6+k(k-1)c_0^{k-2}a_2a_4
                    +\binom{k}{3}c_0^{k-3}a_2^3,
 \qquad a_j=(13/2)_jc_j.
\]
The accompanying source-cost continuation (GS.41)--(GS.50) proves
the first four monic norms $h_0,h_1,h_2,h_3$ and all three positive
ratios $a_j^{(k)}=h_j/h_{j-1}$ from these same moments. In particular,
with $\nu=\mu_2/\mu_0$ and the original sum moments $M_j(k)$,
\[
 a_1^{(k)}=k\nu,\qquad
 a_2^{(k)}=\frac{\mu_4}{\mu_2}+(2k-3)\nu,\qquad
 a_3^{(k)}=\frac{M_6-M_4^2/M_2}{M_4-M_2^2/M_0}.
\]
The exact polynomial isometry $P(S)\mapsto P(k/2+iu)$ and its
monic phase $i^j$ preserve these norms and prove the original
spectral recurrence $SP_j=P_{j+1}+(k/2)P_j-a_j^{(k)}P_{j-1}$.
For the source determinant corrections, put
\[
 Q_N=\frac{X_{N+2}X_N}{X_{N+1}^2},\qquad
 Q_0=\frac{2\nu}{13},\qquad
 Q_1(k)=Q_0+\frac{\rho}{13k+2},\qquad
 \rho=\frac{\mu_4}{\mu_2}-\frac{43\nu}{13}.
\]
The certified interval for $\rho$ is strictly negative; hence
$0<Q_1(1)\le Q_1(k)<Q_0$ and
$Q_1(k+1)-Q_1(k)=-13\rho/((13k+2)(13k+15))>0$.
These are exact statements at every tensor degree for the specified
source polynomial degrees. The analytic seed has the original zero
arithmetic quotient $\mathbb C[S]/(1)$ throughout.

\paragraph{Result R40: finite error propagation to the actual consecutive ratio.}
The exact polynomial coefficient power and test expansion in
(GMT.12)--(GMT.17) propagate finite coefficient enclosures into
Hermitian source envelopes. The full boundary-resolvent formula
(GMT.19)--(GMT.20) carries those actual forms to their least lifts,
including the signed quadratic loss and its original theta primitive.
Equations (GMT.29a)--(GMT.29b) also certify every coefficient of
the least lift and the full representative loss
$\widehat R^*H\widehat R-G(H)$ by an explicit positive
$q$-dimensional rational upper form.
The positive finite envelopes induce quotient-dimensional determinant
enclosures by minimizing on the same affine space of representatives.
For $N\geq q$, retain the exact two-degree source block
\[
 H_+=\begin{pmatrix}H_-&C_D\\C_D^*&D_D\end{pmatrix},\quad
 W=D_D-C_D^*H_-^{-1}C_D,\quad
 F_D=J_D-J_-H_-^{-1}C_D,\quad Z_D=F_D^*G_-F_D.
\]
Equations (GMT.22)--(GMT.27) prove, in the original coordinates,
\[
 \frac{V_{N-1}}{V_{N+1}}=\frac{\det(W+Z_D)}{\det W},\qquad
 a_{N+1}=\frac{\det W}{W_{11}^2}.
\]
Applying the same maps to the Gamma reference gives the exact
relative correction $T_{N-1}/T_{N+1}$ as the quotient of these
two positive ratios. The fully proved auxiliary bounded multiplier
$\min(1+t^2/4,10001)$ with $\lambda=1$, $k=2$, and the declared
$\chi(S)=(S-1)^2-2$ has $T_1/T_3<1$. Its exact finite tail,
rational envelopes, relation matrices, and source/quotient maps
are (GMT.32)--(GMT.42). Thus positivity and boundedness alone
do not imply relative correction monotonicity. This explicit
counterexample is connected to the arithmetic problem by those
same coefficient, least-lift, and determinant maps; its auxiliary
weight is not assigned the value of an actual zeta packet weight.

\paragraph{Result R41: the complete phase and a strict relative contribution.}
For the actual complex product $a=\prod_i A_{h,\lambda}(t_i)$,
put $a_k=C_\Sigma a$, $b_k=C_\Sigma|a|^2$, and
$d_k=b_k-|a_k|^2$. Equations (GP.6)--(GP.11) prove the exact
unitary transport of the arithmetic sum projection and its fibre
defect. Equations (GP.12)--(GP.13) use the original quartet
tail, with its full constant $25\,2^{2d}/\sqrt\pi$, to prove
$d_k(u)>0$ at every real $u$ for $\lambda=1/4$, $d\geq4$,
and $k\geq2$. The full complex coefficient transform and its
Parseval identity preserve the factors $c_\lambda^k$ and $n!$.
The source matrices and the original quotient satisfy
\begin{gather*}
 M_h=M_{\rm coh}+M_{\rm rel},\qquad
 R_h=R_{\rm coh}-B_\chi K,\\
 K=(B_\chi^*M_hB_\chi)^{-1}B_\chi^*M_{\rm rel}R_{\rm coh},\\
 G_h=G_{\rm coh}+K^*B_\chi^*M_{\rm coh}B_\chi K
                            +R_h^*M_{\rm rel}R_h.
\end{gather*}
For every admitted nonzero quartet quotient the last difference
is positive definite. The empty boundary at $N=q-1$ and zero
quotient at $q=0$ retain their explicit separate values. The
graph-to-source addition, gamma multiplication, inverse Mellin
map, and full jet unit in (GP.24)--(GP.27) reconstruct the
original arithmetic observation and its exact supported zero.

The new finite coefficients, strict fibre contribution, and interval
certificates therefore enter the same original Toda and exterior
quantities in R33--R37. In particular the lower spectral cost and
the upper expression are compared in the same $G_N$, with its
volume-imbalance square, phase square, spectral coupling, and
projection terms retained. The following arithmetic endpoint continuation proves a norm bound
for the original measure and a finite upper comparison at the full
amplified degree. Results R42--R46 below give their complete maps
and the explicitly calculated gap between the resulting upper and
lower volume estimates.


\subsection{Arithmetic endpoint estimates and the full reflected relation map}

\paragraph{Result R42: the actual monic norm and its exact leading line.}
The complete arithmetic endpoint source and its analytic audit prove,
for the original density and every integer $n\ge k\ge3$,
\[
 \left(\omega_{h,k,2n}/\omega_{h,k,n}\right)^{1/(2n)}
 \le C_h n.
\]
Here the literal constant $C_h$ is (EP.40), and the original lower
convolution envelope is (AU.10), with $B=42+2\deg h$ and the
compact mass $\vartheta_h^{k-3}$ unchanged. The complete
balanced-window proof strengthens the admitted norm comparison to
\[
 \left(\omega_{h,k,n+r}/\omega_{h,k,n}\right)^{1/(2r)}
 \le C_h^{\rm bal} n,
 \qquad n\ge k\ge3,\quad n/2\le r\le n,
\]
with every factor of $C_h^{\rm bal}$ stated in (EP.50b).
The analytic proof derives the local zeta mass from its exact ellipse,
the entire identity $g=hv_h$ at the selected centres, the three-factor
positive convolution, and the monic Legendre norm. It includes the
two Laplace values and their factor two before the final root.

Equations (AL.3)--(AL.10) prove the endpoint and metric of the
underlying degree map. Multiplication by the original monic $\chi$
gives $\mathcal L_{n-q}\simeq\mathcal L_n$ for $n\ge q$, with
squared operator norm $\omega_n/\omega_{n-q}$ and adjoint
coefficient equal to that same positive ratio. At $n=q-1$ its
source is zero and its target is the original one-dimensional
leading line. For $n\ge q$, its representative has the complete Pythagorean
decomposition
\[
 \|\mathcal V_{h,k}(\chi p_{n-q})\|^2
 =\omega_n+\|\mathcal V_{h,k}(\chi p_{n-q}-p_n)\|^2.
\]
The correction reduces to $-[p_n]_\chi$ and the relation reduces
to zero. Thus the degree-line norm, the original source cost and
the supported arithmetic zero are connected by the explicit
commutative diagrams of the leading-line chapter.

\paragraph{Result R43: exact products, their optimizer, and full window transport.}
Keep $n\ge q$, $r\ge1$, $m=n+r$ and $d_j=V_{n+j}/V_{n+j-1}$.
The full radius calculation in (EP.1)--(EP.11), including the
original complex phase $b_N^*G_Nb_{N+1}=i\omega_N\phi_N$, gives
\[
 \prod_{j=0}^{r-1}(\epsilon_{n+j}^2+\phi_{n+j}^2)
 =\frac{\omega_m}{\omega_n}\frac{V_n}{V_m}
    (1-d_0)(1-d_r)\prod_{j=1}^{r-1}(1-d_j)^2.
\]
For the combined budget
$\mathcal B=-\log(d_0d_r\prod_{j=1}^{r-1}d_j^2)>0$,
the unique maximizing contraction sequence for this scalar product is
\[
 d_0=(1+2x)^{-1},\qquad d_r=1-2x,\qquad
 d_j=(1-x)/(1+x)\quad(1\le j<r),
\]
where $0<x<1/2$ solves
$\mathcal B=2\operatorname{artanh}(2x)+4(r-1)\operatorname{artanh}x$.
The proof in (EP.22)--(EP.31) establishes existence, uniqueness,
the exact maximum, its strict improvement over the earlier
hyperbolic-sine bound when $r>1$, and rational root certificates.
The fixed-four-endpoint optimum, the phase polynomial and the
quantitative concavity loss are separately proved in (EP.15)--(EP.21).
No assertion that the scalar optimizer is realized by an arithmetic
source is needed for the proved bound on the original evaluation map.

The full Hermitian window chapter constructs its Schur matrices
$W_j,F_j,Z_j$ and their determinant ratios from the same source
coefficients. Its finite metric transport retains the eigenvalue
spread and the exact hyperbolic coordinate for the four-volume
logarithm. Equations (EP.32)--(EP.38) give the literal evaluation
map from those matrices to every $d_j$, the budget and the product.
With $H_+,p,e$ defined in (EP.34), the complete cancellation is
\[
 \frac{\omega_m}{\omega_n}\frac{V_n}{V_m}
 =\frac{H_+}{pe}
 =\frac{\mathfrak B_{m-q+1}\mathfrak D_n}
        {\mathfrak D_m\mathfrak B_{n-q+1}}.
\]
The confluent transfer proof (EP.38a)--(EP.38k) constructs the
raw-derivative matrices for every node and order, proves their
invertibility, and proves
$\det\mathsf F_a/\det\mathsf V=\mathfrak B_a/\mathfrak D_a$.
The sign of the column shift, all derivative factorials and every
power of $i$ remain in that construction. It also proves the
theta derivative giving the phase and the exact Gamma-reference
ratio $\det\mathsf F_a/\det\mathsf F_a^\Gamma=Y_a/X_a$.

\paragraph{Result R44: the necessary central and four-volume costs.}
For an actual off-line quartet with full common multiplicity $m_0$,
the retained degree and allowance are
\begin{gather*}
 q_k=[1+k(m_0-1)](k+1)^2,\\
 L_{h,k}=2\delta[1+k(m_0-1)](k+1)
                \left\lfloor (k+1)^2/4\right\rfloor.
\end{gather*}
The complete local tensor calculation and Hermitian compression
proof in (EP.41) give $\epsilon_{h,k,N}\ge L_{h,k}$ and
$L_{h,k}/q_k\ge\delta k/2$. At $n=q_k,r=q_k-1$, write
$C_k=\log(V_{q_k}/V_{2q_k-1})$. Combining the just-proved
balanced norm estimate with the exact product yields
\[
 C_k\ge2(q_k-1)\log\frac{L_{h,k}}{C_h^{\rm bal}q_k}
\]
and the stronger full phase and contraction penalties of (EP.50c).
The exact four-volume identity is
\[
 \log\frac{V_{q_k-1}V_{q_k}}{V_{2q_k-1}V_{2q_k}}
 =2C_k+\log\frac{V_{q_k-1}}{V_{q_k}}
          +\log\frac{V_{2q_k-1}}{V_{2q_k}}.
\]
Both additional logarithms are nonnegative by minimization over
nested lift spaces. Substitution of $L_{h,k}/q_k\ge\delta k/2$,
division by $q_k\log k$, and $(q_k-1)/q_k\to1$ prove the
necessary four-volume lower limit $4$ in (EP.50e).
Equations (EP.44)--(EP.48) construct the original relation layer
and the exact change of least lift. Its central Gram has largest
relative eigenvalue at least
$\delta^2k^2/(4(C_h^{\rm bal})^2)-1$ by (EP.50f).
The cochain primitive and supported quotient image are written
there for the same representative.

\paragraph{Result R45: an actual arithmetic upper comparison and its calculated gap.}
The new arithmetic upper-volume chapter constructs the positive
measure
\[
 d\nu_k(u)=a_k e^{-(\pi/2)|u|}(k-2+|u|)^{-B}\,du,
 \qquad a_k=c_h\vartheta_h^{k-3}e^{-(\pi/2)(k-3)},
\]
which is dominated by the original $m_{h,k}(u)\,du$ for $k\ge3$.
The identity on polynomials is contractive and commutes with every
original raw jet, remainder and spectral multiplication map.
Its moments are the complete positive Gamma-mixture integrals
(AU.12). Writing $\mathsf K_N$ for the original raw-jet kernel,
minimization and positive inversion prove
\[
 0\le C_k\le
 \log\det\mathsf K_{2q_k-1}^{\nu}
       -\log\det\mathsf K_{q_k}.
\]
The actual arithmetic moments needed by this bound stop at degree
$2q_k$; the comparator moments through $4q_k-2$ have the explicit
positive integral in (AU.12). The exact finite Gamma coefficient
map (AU.16) gives the same cutoff on $c_{h,j}$.
The two-Laplace-value majorant, positive Cauchy--Binet minor sum,
and interval Legendre jet kernel in (AU.17)--(AU.23) supply
further proved upper expressions for this same $C_k$.

The full scalar proof (AU.24)--(AU.33) gives
\begin{align*}
 C_k\le\mathcal U_k={}&(\alpha D+\log256)q_k^2\\
 &+\left(\log(XK)+\alpha+\frac{16R_0}{3D}\right)kq_k\\
 &+Bq_k\log(1+(D+1)q_k)+q_k\log C_{\rm env},
\end{align*}
where $\alpha=\pi/2$, $R_0=\sqrt{\delta^2+\gamma^2}$,
$D=\max(1,2/b,2R_0)$, $0<b<\pi/2$,
$X=\max(1,M_h(b),M_h(-b))$, $K=\max(1,\vartheta_h^{-1})$,
and $C_{\rm env}=\max(1,2/(3c_hD))$.
All constants and the literal source mass are retained. For
$\mathcal L_k=2(q_k-1)\log(\delta k/(2C_h^{\rm bal}))$,
(AU.32) computes every term of $\mathcal U_k-\mathcal L_k$,
and (AU.33) proves that its quotient by $q_k^2$ tends to
$\alpha D+\log256>0$. This upper estimate leaves the stated
gap to the quartet lower estimate.

The complete auxiliary Lebesgue calculation proves
\[
 \left|\log(V_q^{\rm aux}/V_{2q-1}^{\rm aux})
       -\left(\tfrac92\log3-6\log2\right)q^2\right|
 \le12q\log(8q)
\]
for every even $q\ge2$. Its full-origin-jet quotient is connected
to the original grid by the free monic family
$\mathbb C[t,u]/\prod_{a,b}(u-t\zeta_{ab})^\ell$ in
(AU.46)--(AU.48). The original full jets at $t\ne0$, the merged
jet at zero and the source-dependent kernel derivative are all
explicit. The displayed auxiliary estimate belongs to its
Lebesgue source; the exact derivative gives the change along the
family, without assigning the same asymptotic to the arithmetic source.

\paragraph{Result R46: reflected norm, conormal algebra and their complete kernels.}
For any original monic $\chi$ of degree $q\ge1$, set
$\chi^\dagger=(-1)^q\overline\chi(k-S)$,
$N_\chi=\chi\chi^\dagger$ and $g=\gcd(\chi,\chi^\dagger)$.
The reflected relation chapter proves the algebra isomorphism
\[
 \mathbb C[S]/(\chi^2\chi^\dagger/g)
 \simeq
 \mathbb C[S]/(\chi^2)
 \mathbin{\times}_{\mathbb C[S]/(\chi g)}
 \mathbb C[S]/(\chi\chi^\dagger),
\]
with its exact sequence, inverse residue construction and original
$M_S$ action. The norm map retains
$\chi(k/2+iu)=i^q\psi(u)$ and
$N_\chi(k/2+iu)=i^{2q}\psi(u)\overline\psi(u)$.
Thus the full union of reflected and original roots, with added
orders at overlaps, is the precise general norm-transfer domain.

The kernel coefficient modules have annihilators $(\chi)$ and
$(\chi^\dagger)$. Equations (RF.19)--(RF.21) prove their
conjugate-linear bijection
$F([\chi f]_{\chi^2})=[\chi f^\star]_{N_\chi}$ with inverse,
$FM_S=(kI-M_S)F$, and the exact transported norm-kernel product.
If $a_\lambda,b_\lambda$ are the respective root orders,
\[
 \dim I_N^t=\sum_\lambda
            \max(b_\lambda-(t-1)a_\lambda,0),\qquad t\ge1.
\]
Equations (RF.22)--(RF.26) prove the exact nilpotence index, the
persistent reflected-only CRT algebra, and the fully worked mixed
example of kernel dimensions $3,1,0$. Projection sends each
specified kernel vector, including the persistent sector's
idempotent when present, to supported zero in its original fibre.
The same lifted map fixes the external point $\tau$.



\paragraph{Result R47: restriction spectra, actual arithmetic gaps and finite certificates.}
For every $q-1\le i\le j$, retain the original quotient kernels
$K_N$, metrics $G_N=K_N^{-1}$ and minimum source lifts $R_N$.
In the actual monic orthogonal coordinates, let $L_{i,j}$ pad source
coefficients and let $P_{i,j}$ be the orthogonal endomorphism
projection onto its range in $\mathcal P_j$. The complete
restriction source and the arithmetic join prove
\[
 T_{i,j}=K_iG_j:(E,G_j)\longrightarrow(E,G_i),\qquad
 L_{i,j}R_iT_{i,j}=P_{i,j}R_j.
\]
Let $I_{i,j}:(E,G_i)\to(E,G_j)$ denote the original identity
transport, and let $A_i,A_j$ denote the same original action $A$
on these two specified metric spaces. The return is its metric
adjoint. The actual endomorphism
$U_{i,j}=T_{i,j}I_{i,j}$ of $(E,G_i)$ has coordinate matrix
$K_iG_j$. Its eigenvalues $0<g_a\le1$
are the squared singular values of the source restriction between
the two canonical minimum-lift ranges. For the orthogonal projections
$\Pi_i,\Pi_j$ onto those ranges in the same original source, the
join constructs both restricted-product similarities and proves
\[
 \operatorname{Tr}_{\mathcal P_j}|\Pi_i-\Pi_j|^{2m}
   =2\operatorname{Tr}_E\bigl((I-U_{i,j})^m\bigr),
 \qquad m\ge1.
\]
Its explicit two-dimensional principal-angle blocks retain all
common-range and ambient-zero components.

The original arithmetic action $A$ has the exact reverse-map defect
\[
 A_iT_{i,j}-T_{i,j}A_j
  =\mathsf H_iT_{i,j}-T_{i,j}\mathsf H_j,
 \qquad
 \mathsf H_N=K_N(A^*G_N+G_NA-kG_N).
\]
The change of representatives is the complete relation
\[
 \mathfrak d_{i,j}=L_{i,j}R_i-R_jI_{i,j}
   =L_{i,j}R_i(I-U_{i,j})-(I-P_{i,j})R_jI_{i,j},
 \qquad
 \mathfrak d_{i,j}^*O_j\mathfrak d_{i,j}=G_i-G_j.
\]
Here $\mathfrak d_{i,j}:(E,G_i)\to\mathcal P_j$ has the
original theta image. Both terms on the right have the same
remainder $I_{i,j}(I-U_{i,j}):(E,G_i)\to(E,G_j)$:
indeed $B_jL_{i,j}R_i=I_{i,j}$ and
$B_jP_{i,j}R_jI_{i,j}=I_{i,j}U_{i,j}$. The full division primitive and theta image of their
difference are supplied in the join. Reduction sends its labelled
relation vector to the supported zero at that receiving label;
the lifted map sends external $\tau$ to $\tau$.

The arithmetic remainder estimate is extended to every admitted
degree with its entire monic quotient formula. In the original
notation of the upper-comparison chapter, it gives
\[
 G_{q-1}\preceq C_NG_N,\qquad
 C_N=\frac{U_k(T)A_N^2B_q(r)^2}{T a_k(T)},\qquad
 N\ge q-1,
\]
with $T=Dq$, $r=kR_0/T<1$, the unchanged original density mass,
and all constants defined in the complete proof. Thus the actual
endpoint gaps $C_{2q-1}^{-1}$ and $C_{2q}^{-1}$ are proved
arithmetic inequalities. The exact ratio satisfies
$C_{2q}/C_{2q-1}\le147/5<32$, so the same construction gives
$\mathcal B_{h,k}\le2\mathcal U_k+q\log32$ with
$\mathcal U_k$ from R45.

The finite matrix construction improves these directional bounds.
For rational Hermitian enclosures of the actual kernels,
\[
 0<L_i\preceq K_i\preceq U_i,\qquad
 0<L_j\preceq K_j\preceq U_j,
\]
define
\[
 g=\bigl(\operatorname{Tr}(L_i^{-1}U_j)-q+1\bigr)^{-1},
 \quad r_*=1-g,\quad
 H_*=I-L_iU_j^{-1},\quad t_m^*=\operatorname{Tr}(H_*^m),
\]
\[
 d=\operatorname{Tr}(U_iL_j^{-1})
     -\operatorname{Tr}(L_iU_j^{-1}).
\]
The complete ordered generalized-eigenvalue proof gives
$0<g\le1$, $d\ge0$ and the rational intervals
\[
 \max\{0,t_m^*-m r_*^{m-1}d\}
 \le\operatorname{Tr}\bigl((I-U_{i,j})^m\bigr)\le t_m^*.
\]
The exact quotient-first triangular basis constructs these brackets
from certified source and relation-coefficient intervals, retaining
the actual polynomial rather than replacing it by an interval centre.
For $r_*>0$, summing the moment intervals and bounding the positive
logarithmic tail gives a certificate for $\log(V_i/V_j)$ with width
at most
\[
 \frac{d(1-r_*^p)}g+
 \frac{q r_*^{p+1}}{(p+1)g}.
\]
The source's sharper coefficient and a rational logarithm enclosure
give the further proved refinement. At $g=1$ the loss is exactly
zero. Entry refinement followed by increasing $p$ proves convergence
at each fixed pair with every error term displayed.

The balanced central product of R44 now has a consequence for each
of the two original endpoint pairs. Put
$X=L_{h,k}/(C_h^{\rm bal}q)$ and retain the full nonnegative
phase/contraction penalty $P_{\rm mid}$ from the central product.
Since $\operatorname{rank}(K_{2q-1}-K_q)\le q-1$, at least
one central return eigenvalue is exactly one. The remaining
$q-1$ eigenvalues have product at most
$X^{-2(q-1)}e^{-P_{\rm mid}}$. The ordered Rayleigh-quotient
comparison with each endpoint therefore proves
\[
 \min\operatorname{Spec}U_{q-1,2q-1},\quad
 \min\operatorname{Spec}U_{q,2q}
 \ \le\min\left\{1,X^{-2}e^{-P_{\rm mid}/(q-1)}\right\}
 \le\min\left\{1,\frac{4(C_h^{\rm bal})^2}{\delta^2k^2}\right\}.
\]
Every such eigenvalue is the squared norm fraction of the actual
higher-degree minimum representative retained by the lower-degree
source projection. The inverse return recovers its original
coordinate, with precisely the relation correction proved above.



The whole inverse spectrum has the corresponding finite lower bound.
To retain endomorphism types, write
$U_{i,j}=T_{i,j}I_{i,j}$ on $(E,G_i)$, where $I_{i,j}$ is the
original identity transport. Its coordinate matrix is $K_iG_j$.
For every real $s>0$, each endpoint pair separately obeys
\[
 \operatorname{Tr}U_{i,j}^{-s}
 \ge\max\left\{q,\,
 1+(q-1)X^{2s}e^{sP_{\rm mid}/(q-1)}\right\}.
\]
Indeed the central trace is one plus the inverse-power sum of its
remaining $q-1$ eigenvalues. The tangent inequality $e^y\ge1+y$
at their logarithmic mean bounds this sum from below by their
geometric mean, and the ordered comparison transfers every eigenvalue.
If the exact central update rank is
$\rho=\operatorname{rank}(K_{2q-1}-K_q)>0$, put
$A_+=\max\{0,2(q-1)\log X+P_{\rm mid}\}$. Retaining all
$q-\rho$ unit returns gives the stronger bound
\[
 \operatorname{Tr}U_{i,j}^{-s}
 \ge q-\rho+\rho e^{sA_+/\rho}.
\]
At $\rho=0$ the central trace is exactly $q$ and both endpoint
traces are at least $q$. The complete proof and rank comparison are
(ERJ.42)--(ERJ.42a).
Under the canonical isometries the source restriction is the
bijection $\mathcal R=\Pi_i|_{\ell_jE}$ between the two canonical
images. Its inverse singular-value-power sum is exactly
$\sum_a\sigma_a(\mathcal R^{-1})^{2s}
=\operatorname{Tr}U_{i,j}^{-s}$ by (ERJ.11a).
The ambient zero eigenspaces of the full projection products remain
outside this restricted inverse and are retained by (ERJ.9).


\paragraph{Result R48: the complete arithmetic weight optimizer and a stronger moment family.}
For the original relation $\chi$, coordinate $S=k/2+iu$ and full
raw-jet matrix $E_q$, put $\psi(u)=i^{-q}\chi(k/2+iu)$,
$Z=E_{q-1}$ and $\beta_j=b^{2j}/(2j)!$ for $0<b<\pi/2$.
The weight-optimization chapter proves the signed cofactor map
\[
 \det E_q[:,\{0,\ldots,q\}\setminus\{j\}]
   =(-1)^{q-j}\psi_j\det Z,
 \qquad \psi_j=\frac{i^{j-q}}{j!}\chi^{(j)}(k/2).
\]
In particular every vanishing coefficient has its corresponding
zero minor, with the full jet space retained. For auxiliary weights
$\tau_j\ge0$ of sum one, set
\[
 \eta_j=\frac{(2j)!}{(j!)^2b^{2j}}|\chi^{(j)}(k/2)|^2,
 \quad D_\eta(\tau)=\sum_{j=0}^q\eta_j\prod_{l\ne j}\tau_l,
 \quad d(\eta)=\max_{\tau\ge0,\,\sum\tau_j=1}D_\eta(\tau).
\]
Equations (OW.4)--(OW.13) prove the precise determinant identity
\[
 \det(E_q\operatorname{diag}(\tau_j\beta_j)E_q^*)
 =|\det Z|^2\Bigl(\prod_{j=0}^q\beta_j\Bigr)D_\eta(\tau)
\]
and its original phase-preserving congruence. The chapter proves
every maximizer: a largest coefficient at least the sum of the
others gives the corresponding zero weight and the remaining
weights $1/q$; exactly two equal positive coefficients give the
explicit segment of maximizers. Otherwise the maximizer is unique
and interior. Its zero-coefficient weights are $1/q$, and its
positive-coefficient weights are $(1-y_j)/q$, where
\[
 \sum_{\eta_j>0}y_j=1,\qquad
 \eta_j=c_*y_j(1-y_j),\qquad 0<y_j<1.
\]
The two monotone scalar constructions (OW.16)--(OW.19) determine
the unique solution, including the branch with one large root.
Equations (OW.21)--(OW.22a) prove the rational three-coefficient
formula, coefficient-error bounds and terminating rational value
certificates. The linear case $q=1$ and all boundary cases are
included in the same complete proof.

With the literal original $M_k=M_h(b)^k+M_h(-b)^k$ and comparator
volume $V_{2q-1}^{\nu}$ from R45, the resulting arithmetic bound is
\[
 C_k\le\mathcal U_{\rm opt}
   =q\log M_k-\log\!\Bigl(\prod_{j=0}^q\beta_j\Bigr)
      -\log d(\eta)-\log V_{2q-1}^{\nu}.
\]
The limit of the original positive raw-kernel inequality proves
this bound at boundary optima; no reciprocal zero weight is used.
For $\mathcal U_{\rm eq}$ obtained from the same inequality with
equal weights, (OW.23)--(OW.29) prove
\[
 0\le\mathcal U_{\rm eq}-\mathcal U_{\rm opt}
   =\log\frac{d(\eta)}{\sum_j\eta_j/(q+1)^q}
   \le q\log(1+1/q)<1.
\]
Both non-strict inequalities are strict for the admitted complete
quartet: its parity gives a zero coefficient, while its nonzero
roots and monicity give at least two nonzero coefficients. Thus
the finite optimization within this specified family is solved,
and its improvement divided by $q_k\log k$ or $q_k^2$ tends to zero.

The continuation (OW.30)--(OW.40) proves a stronger source inequality.
Write $\mu_{2j}=\int u^{2j}m_{h,k}(u)\,du$, retaining
$\mu_0=\mu_h^k$, and set $S_q(b)=\sum_{j=0}^q\beta_j\mu_{2j}$.
Pointwise Cauchy--Schwarz and the full positive even exponential
series, integrated against the original source, give
\[
 H_q\preceq S_q(b)\operatorname{diag}(\beta_j^{-1})
   \prec\frac{M_k}{2}\operatorname{diag}(\beta_j^{-1}),
 \qquad S_q(b)<M_k/2.
\]
The corresponding upper expression $\mathcal U_{\rm cosh}$ improves
the earlier $\mathcal U_{\rm eq}$ by exactly
$q\log(2(q+1))$. Using $S_q(b)$ gives the further strict
improvement $q\log(M_k/(2S_q(b)))$. Every constant follows from
the integrated finite sum and both original Laplace values.

For every positive auxiliary list $r_j$, the same calculation proves
\[
 H_q\preceq\left(\sum_{j=0}^q r_j\mu_{2j}\right)
                \operatorname{diag}(r_j^{-1}).
\]
The exact positive-cone map
$t_j=r_j\mu_{2j}/\sum_l r_l\mu_{2l}$, with inverse
$r_j=s t_j/\mu_{2j}$, transforms its determinant optimization
to the already solved polynomial $D_\kappa$, where
$\kappa_j=|\psi_j|^2\mu_{2j}$. Thus the entire stronger
diagonal family has the proved optimum
\[
 C_k\le\mathcal U_{\rm mom}
   =\sum_{j=0}^q\log\mu_{2j}-\log d(\kappa)
                   -\log V_{2q-1}^{\nu}.
\]
All $q+1$ moment factors, every zero coefficient and the boundary
raw-kernel limit are retained. The ray $r_j=\beta_j$ maps to
$t_j^{\beta}=\beta_j\mu_{2j}/S_q(b)$, and the exact total gain is
\[
 \mathcal U_{\rm opt}-\mathcal U_{\rm mom}
 =q\log(2(q+1))-\mathfrak g(\eta)
   +q\log\frac{M_k}{2S_q(b)}
   +\log\frac{d(\kappa)}{D_\kappa(t^{\beta})}.
\]
The last term is nonnegative and is evaluated by the same complete
optimizer. The removed term $q_k\log(2(q_k+1))$, divided by
$q_k\log k$, tends to $2$ for a simple quartet and to $3$ for
fixed common multiplicity greater than one.

Finally the original quotient-first source Gram proves
$V_q=\det H_q/(\boldsymbol\psi^*H_q\boldsymbol\psi)$.
Equation (OW.40) therefore calculates the exact remaining comparison
loss against the full low-degree kernel:
\[
 \mathcal U_{\rm mom}
       -\bigl(\log V_q-\log V_{2q-1}^{\nu}\bigr)
 =\sum_{j=0}^q\log\mu_{2j}-\log d(\kappa)-\log\det H_q
      +\log(\boldsymbol\psi^*H_q\boldsymbol\psi)\ge0.
\]
These identities retain all actual moments through degree $2q$.
They determine the finite improvement and its precise comparison
map without assigning an unproved growth rate to the last loss or
the comparator determinant. The original source, support label and
quotient action remain connected through the same cofactor and
kernel maps.


\paragraph{Result R49: the first-degree phase product and each full block budget.}
The complete first-window chapter joins the original first-degree
radius to the comparable-window theorem of PR~25, retaining the
nonnegative generalized-eigenvalue correction. At $N=q-1$, expand
$\chi=p_q+\sum_{j<q}\gamma_jp_j$ and keep
$\nu_0=\|\chi\|^2=\omega_q+\sum_{j<q}|\gamma_j|^2\omega_j$.
The rank-one update and the original phase identity give
\[
 \frac{V_q}{V_{q-1}}=\frac{\omega_q}{\nu_0},\qquad
 \gamma_{q-1}=-i\phi_{q-1},\qquad
 \epsilon_{q-1}^2+\phi_{q-1}^2
       =\frac{\nu_0-\omega_q}{\omega_{q-1}}.
\]
Equations (CJ.3)--(CJ.7) prove each equality in the actual source.
For $\Lambda_N=\omega_N/V_N$, the first step consequently has
exactly one contraction:
\[
 (\epsilon_{q-1}^2+\phi_{q-1}^2)\Lambda_{q-1}
       =\Lambda_q(1-\delta_q).
\]
Multiplication by the regular consecutive identities gives the full
first-window formula (CJ.9), for every $j\ge q$:
\[
 \prod_{N=q-1}^{j-1}(\epsilon_N^2+\phi_N^2)
 =\frac{\omega_jV_{q-1}}{\omega_{q-1}V_j}
       (1-\delta_j)\prod_{N=q}^{j-1}(1-\delta_N)^2.
\]
The original source contains no quotient inverse below degree $q-1$.
The complete two-sided monic norm envelope (CJ.15)--(CJ.18)
supplies the explicit $C_h^{\rm win}$ for $n/2\le r\le2n$.
For $D_h=\delta/(2C_h^{\rm win})$ and each original block
$(i,j)=(q-1,2q-1)$ or $(q,2q)$, it proves the finite bound
\[
 \log(V_i/V_j)\ge2q\log(D_hk)+\mathcal P_{i,j}
      +\sum_{N=i}^{j-1}\log(1+\phi_N^2/L_{h,k}^2).
\]
The penalties $\mathcal P_{i,j}$ are written completely in (CJ.13),
including the single terminal factor for the first block and both
boundary factors for the regular block. All terms added to the
per-block $2q$ budget are nonnegative.

The original relation map in (CJ.20)--(CJ.21) has Gram
$\Omega+F^*G_iF$ and is an isomorphism onto the whole newly
admitted relation layer. At zero tilt and even tensor degree,
the exact quartet's first increment matrix has a kernel by the
full parity calculation. Every zero generalized eigenvalue stays
in the $q$-dimensional cost spectrum, contributing the factor one
to the determinant. For all $s>0$, (CJ.22) also proves
\[
 \operatorname{Tr}Q_{i,j}^{-s}
       \ge\max\{q,q(D_hk)^{2s}\},
 \qquad Q_{i,j}=T_{i,j}I_{i,j}\text{ on }(E,G_i).
\]
These are traces of the same original metric endomorphism. The
full cost, source relation, and its supported-zero quotient image
remain connected by the maps in that chapter.


\paragraph{Result R50: the two relation traces and their actual comparison source.}
The complete relation-moment chapter proves the sharp finite determinant
majorant from the two traces of the original positive cost operator
\(Z=(K_j-K_i)G_i\), keeping its full dimension and every zero eigenvalue.
Writing \(t_1=\operatorname{Tr}Z\), \(t_2=\operatorname{Tr}Z^2\), for
\(q\ge2\) define
\[
 d=\sqrt{\frac{qt_2-t_1^2}{q-1}},\qquad
 a=\frac{t_1+(q-1)d}{q},\qquad b=\frac{t_1-d}{q}.
\]
The complete rational Hermite remainder in that chapter gives
\[
 \log(V_i/V_j)=\log\det(I+Z)
 \le\log(1+a)+(q-1)\log(1+b).
\]
For dimension one the exact expression is \(\log(1+t_1)\), and for
dimension zero it is zero. The proof identifies both traces with
the original relation map and its exterior square, including every
cross pairing. For the arithmetic upper kernel \(U_j\), the positive
space \(\mathcal C_+=(E^\vee/\ker(U_j-K_i),U_j-K_i)\) and its map
\(J_+[\alpha]=(U_j-K_i)\alpha\) realize the comparison cost
\(Z_+=(U_j-K_i)G_i\). The constructed contraction from the original
relation source has the complete positive defect
\[
 J_+(I-CC^\dagger)J_+^\dagger=(U_j-K_j)G_i.
\]
This retains the precise map and the excess in the bound, together
with the original arithmetic action defect and all unit return factors.
The relation-moment chapter contains the full definitions and proofs.

\paragraph{Result R51: the full Deligne comparison and the constituent coupling.}
The detailed Deligne sidebar proves the image, kernel and cone maps
for the same compact-to-ordinary comparison, and the complementary
degree duality with its actual boundary line. For every specified
invariant \(F\subseteq E\) and divisor depth \(m\ge1\), its original
curve has the logarithmic complex
\[
 \mathscr K_{F,m}=
 [F\mathcal O+E\mathcal O(-m\infty)
   \xrightarrow d(F\mathcal O+E\mathcal O(-m\infty))
                       \otimes\Omega^1(\log\infty)].
\]
Its complete cohomology is \(F\) in degree zero,
\((E/F)\otimes\mathbb T\) in degree two and zero in degree one,
with \(\mathbb T=H^1(\Omega^1)\) and the retained trace
\(\operatorname{tr}[dz/z]=1\). The higher-depth quotient contracts
by division by the displayed nonzero integer coefficients; its
dual retains all meromorphic powers and the twist by
\(\mathcal O(-\infty)\). Each split lift distinguishes the represented
kernel from external absence through their disjoint inclusions into
the full kernel carrier.

The same original metric, \(G(t)=G_j+t(G_i-G_j)\), gives the invariant
quotient metric and coupling map proved in (DS53)--(DS65). With those
specified maps, its endpoint identity is
\[
 \log(V_i/V_j)=
 \log\frac{\det G_F(1)}{\det G_F(0)}
 +\log\frac{\det Q_+}{\det Q_0}
 -\sum_{r=1}^{\infty}\frac{\operatorname{Tr}(\Xi_F^r)}r.
\]
The series has nonnegative terms before subtraction, and the sidebar
proves its finite remainder. It also proves the exact original theta
source defect
\(dK(z)H_t+r(z)Ia_t\), so the theta boundary and the retained
constituent pass through their specified successive quotients.
The two positive determinants can use R50 on their respective
compressed matrices; the coupling subtraction remains in the same
upper estimate.

\paragraph{Result R52: a degree-one exponential family with exact periods.}
For the original cyclic polynomial \(\chi\) and its specified primitive
\(\Phi\), the delivered exponential note proves
\[
 \mathcal H=\operatorname{coker}
 \bigl(u\partial_S+\chi(S)-t\bigr),\quad
 \mathcal H\cong\mathbb C[u,t]^q,\quad
 \mathcal H/(u,t)\mathcal H\cong E,
 \quad(-u\nabla_t)\bmod(u,t)=A.
\]
Its coefficient-module division, original theta observation,
finite-rank source deformation, infinity chart and dual residue
pairing are all retained with complete proofs in the exponential
note and sidebar. The new determinant proof (XD1)--(XD26) computes
the entire period determinant for the stated oriented rays:
\[
 \det\Pi(u,c,t)=\det\Pi^{\mathrm{mon}}(u)
       \exp\!\left(\frac{\operatorname{Tr}\Phi_t(A(t))}{u}\right).
\]
All Gamma factors and monomial phases are given in (XD16)--(XD18).
The proof uses the actual twisted-reduction matrices, keeps their
off-diagonal corrections, and proves that their trace corrections
vanish. It establishes nonvanishing across every repeated-root
parameter. On the original canonical metric it gives
\[
 \log\det B_N^{\mathrm{per}}=
 \log C_q(u)+2\operatorname{Re}
       \frac{\operatorname{Tr}\Phi_t(A(t))}{u}-\log V_N.
\]
Subtracting at the same parameters retains the exact original
endpoint cost. The translation companion (DT1)--(DT13) further
proves the complete maps under \(\chi_a(S)=\chi(S-a)\): the
coefficient action shifts by \(aI\), the period gains the scalar
\(\exp[(-\Phi(-a)-ta)/u]\) and the full Pascal matrix, and the
finite-field realization tensors the corresponding constant
Artin--Schreier character line. The original arithmetic fibre
remains \(a=0\). These explicit maps preserve the canonical metric
comparison and the full inherited endpoint inequalities.


\paragraph{Result R53: the residue trace and the complete collision map.}
For the unchanged monic polynomial \(H(S)=\chi(S)-t\), the complete
SGA trace chapter constructs the perfect residue functional
\(\lambda_t[P]=[S^{q-1}]\operatorname{rem}_H P\) and the actual
diagonal coevaluation
\[
 \mathcal C_H(X,Y)=\frac{H(X)-H(Y)}{X-Y},\qquad
 \operatorname{Tr}_{E_t/R[t]}M_P=\lambda_t(\chi'P).
\]
The fraction is a polynomial, its multiplication image is \(\chi'\),
and the chapter proves the inverse pairing identities without
dividing by a discriminant. The original logarithmic source
relation maps to \([\chi'P]_\chi\) before this trace is taken.
Its complete periodic diagonal resolution retains both odd and
positive even degrees.

The new comparison (SC26)--(SC30) proves the precise specialization
of those degrees. Over \(T=\mathbb C[t]\), put
\(\mathscr E=T[S]/(\chi-t)\) and \(M=\mathscr E/\chi'\mathscr E\).
The isomorphism \(\mathscr E\simeq\mathbb C[S]\), with
\(t\mapsto\chi(S)\), makes
\(0\to\mathscr E\xrightarrow{\chi'}\mathscr E\to M\to0\)
a free \(T\)-resolution. For \(\kappa=\mathbb C[t]/(t-t_0)\),
\[
 \operatorname{Tor}_1^T(M,\kappa)
 \xrightarrow{\sim}\operatorname{ann}_{E_{t_0}}(\chi'),
 \qquad [v]_{\chi'}\longmapsto[w]_{\chi-t_0},
 \quad(\chi-t_0)v=\chi'w.
\]
The proof gives both inverse directions and the exact \(S\)-action.
Thus every positive even group of the family is zero, while the
positive even groups at a collision are these explicit Tor modules;
the odd groups are \(M\otimes_T\kappa\). Each length-\(m\) local
factor retains both \(m-1\)-dimensional modules and their nilpotent
actions. The Gamma calculation also evaluates the entire period
Gram determinant as
\[
 \det(\Pi^*\Pi)=(2\pi|u|)^q
  \exp\left(2\operatorname{Re}
          \frac{\operatorname{Tr}\Phi_t(A+t\mathcal R)}u\right),
\]
with the complete phase-containing determinant still given by R52.

\paragraph{Result R54: exact constituent curvature and neighboring volumes.}
For an actual invariant injection \(I:F\hookrightarrow E\),
\(0<p=\dim F<q\), put
\(\ell[P]=[S^{q-1}]\operatorname{rem}_\chi P\),
\(\mathcal R=e_0\ell\), \(L=\ell I\),
\(Y=\Pi I\), \(H=Y^*Y\),
\(N=1-YH^{-1}Y^*\), and \(z=N\Pi e_0\).
The complete proof (SC9)--(SC25) identifies the invariant image
with the unchanged polynomial ideal \((g)/(\chi)\), proves
\(L\ne0\) and \(z(t)\ne0\) for every parameter, and computes
\[
 NY'=-\frac t u zL,\qquad
 \partial_{\bar t}\partial_t\log\det H
    =\frac{|t|^2}{|u|^2}\|z\|^2LH^{-1}L^*.
\]
The normal derivative has rank one at every nonzero parameter;
the normal acceleration has rank one at zero. The curvature
vanishes precisely quadratically at zero and is positive elsewhere.
Its factorization through the original section \(r_Ne_0L\),
quotient \(j_E\), and period observation \(N\Pi\) retains the
source kernel and its supported zero.

For the literal polynomial frames
\(J_g=[g,Sg,\ldots,S^{p-1}g]\), its prefix \(J_-\), and
\(J_+=[J_g,e_0]\), put
\(h_J=\det(J^*\Pi^*\Pi J)\). The full frame map,
Schur-complement and cofactor proof in (SC31)--(SC37) gives
\[
 \partial_{\bar t}\partial_t\log h_g
       =\frac{|t|^2h_+h_-}{|u|^2h_g^2}.
\]
Each volume factors exactly as \(h_J=v_{N,J}b_{N,J}\), where
\(v_{N,J}=\det(J^*G_NJ)\) retains the original theta metric
and \(b_{N,J}\) is its specified period-comparison determinant.
The empty prefix has determinant one. In codimension one,
\(\det J_+=(-1)^{q-1}\), so \(h_+\) is the full explicitly
evaluated period determinant above. These are the exact maps and
weights joining this curvature to the neighboring-volume calculation.

\paragraph{Result R55: the retained Hochschild and critical-line comparisons.}
The complete Connes--Consani comparison proves the actual restricted
factorization
\[
 V\xrightarrow{j}j(V)\xrightarrow{\operatorname{Tr}}\mathscr B,
 \quad \operatorname{Tr}j=\Theta,
 \quad\mathcal E=\tfrac12\mathcal U\Theta,
 \quad (D^2-D)e^{-\pi u^2}=\phi_*.
\]
The half-density map \(\mathcal U F=u^{1/2}F\) is the specified
half-line isometry. The cochain isomorphism uses the complete
two-moment source and the Laplacian's continuous inverse. The
critical-line map \(\Gamma:Q\to Q_{\rm crit}\) has, on each
original finite packet, the exact kernel
\(\ker(\Gamma\sigma_Z)=E_{Z,\rm off}\), including every
generalized block. The critical recovery is the explicit jet
\[
 \Xi_\rho[f]=\sum_{j=0}^{m-1}
      \frac{i^{-j}}{j!}f^{(j)}(\gamma)\epsilon^j,
 \qquad \rho=\tfrac12+i\gamma.
\]
The full note proves the reciprocal multiplier argument, the
recovery composite and the supported image of every killed vector.

The raw circle zero mode is retained as a separate smooth coordinate;
the specified Gaussian average of the original theta seed proves
that this line is in the closed source-generated module. Its exact
local flat-seminorm estimate and the printed large-circle limit
correction remain with that proof. Finally the original-metric
Laplacian identities
\[
 \mathsf L_k^*G-G\mathsf L_k=A^*W_k-W_kA,
 \qquad G\mathsf L_k=W_kA-A^*GA
\]
and their complete numerical-range argument are retained with the
actual metric control parameter. Neither the trace construction,
the finite-kernel calculation, nor the period determinant supplies
an unproved uniform bound for that parameter.


\paragraph{Result R56: the actual critical-comparison kernel through periods and tensors.}
The full proof (PC1)--(PC41) identifies the single-packet polynomial
\(\chi_{h,1}=h\) by its explicit variable isomorphism
\(\iota_1:[P(S)]\mapsto[P(s)]\). The original arithmetic
inclusion is \(\eta_h=M_{\upsilon_h}\iota_1\), with
\(\upsilon_h=j_h(2\xi/h)\) and every inverse jet coefficient
retained. It proves
\[
 J_hr_N=\eta_h,\qquad j_E=\eta_h^{-1}J_h,\qquad
 j_Er_N=I,\qquad \mathsf q r_N=\sigma_h\eta_h.
\]
The critical comparison is consequently the actual weighted map
\(\mathcal C_h=\Gamma\sigma_h\eta_h\); its kernel is the full
off-critical primary summand. Each surviving critical jet is
recovered with the original \(i^{-j}/j!\) factor and the inverse
of the displayed local arithmetic unit.

For every original integer tensor degree \(k\ge1\), retain
\(E_k=\mathbb C[S]/(\chi_{h,k})\), its actual weighted inclusion
\(\eta_k\), and the factorwise completed-projective comparison
\[
 \mathcal C_k=
 \Gamma^{\widehat\otimes k}\sigma_h^{\widehat\otimes k}\eta_k.
\]
Let \(\chi_{c,k}\) be the monic sum polynomial obtained from the
all-critical tuples, with the complete maximum nilpotent length at
each literally equal sum; for no critical factors set it to \(1\).
The full proof constructs the variable isomorphism on the tensor
carriers, preserves the original unit, and gives a continuous jet
retraction proving that the finite image survives completion. It yields
\[
 \boxed{\ker\mathcal C_k=(\chi_{c,k})/(\chi_{h,k}).}
\]
At a sum \(\lambda\) with full multiplicity \(m_\lambda\) and
critical-tuple multiplicity \(r_\lambda\), its exact local kernel
is the image of
\([b]\mapsto[x^{r_\lambda}a_\lambda(x)b]\) from
\(\mathbb C[x]/(x^{m_\lambda-r_\lambda})\), where
\(x=S-\lambda\) and the local unit in \(\chi_{c,k}\) is retained.
This includes the additional jets when different tuples share a sum.

For the period matrix of this same original polynomial,
\(\mathcal C_{k,t}=\mathcal C_k\Pi_k^{-1}\) has kernel
\(\Pi_k\ker\mathcal C_k\). Its exact sum-action defect is
\[
 \mathcal C_{k,t}A_{k,t}^{\rm per}
       -M^{[k]}\mathcal C_{k,t}
   =t\mathcal C_k(e_{0,k})\ell_k\Pi_k^{-1}.
\]
For a proper nonzero kernel it has the proved rank-one restriction,
with kernel \(\Pi_k(\ker\mathcal C_k\cap\ker\ell_k)\).
The curvature and neighboring-minor identities in R54 therefore
measure this actual critical-comparison kernel. The original theta
metric is connected by the complete congruence and derivative
cancellation. All-critical and no-critical packets retain their
different empty and full kernels, even though both endpoint
curvatures vanish.

\paragraph{Result R57: computed Laplacian control and its curvature comparison.}
The full original-metric calculation (LC1)--(LC16) uses
\(A_t=A+tR\), \(R=e_0\ell\), and
\(W_t=A_t^*G_N+G_NA_t-kG_N\).
For \(q>1\), \(\ell e_0=0\), so \(R^2=0\). The two actual
\(G_N\)-orthogonal columns \(e_0,G_N^{-1}\ell^*\) give
the exact additional control spectrum
\[
 \operatorname{spec}_{\ne0}
     (G_N^{-1}(W_t-W_0))
       =\{\,|t|\sigma_N,-|t|\sigma_N\,\},\qquad
 \sigma_N^2=(e_0^*G_Ne_0)(\ell G_N^{-1}\ell^*),
\]
when \(t\ne0\). For \(q=1\) the same difference is exactly
\(2\operatorname{Re}t\) times identity. For the retained
reflection-stable packet, the complete earlier source-minor calculation evaluates
\(\epsilon_N^2=\alpha_{N+1}\Delta_N/\det K_N\), so the actual
quantity
\[
 E_N(t)=\epsilon_N+
 \begin{cases}|t|\sigma_N,&q>1,\\
               2|\operatorname{Re}t|,&q=1
 \end{cases}
\]
bounds the control on the same metric. With this computed value,
every point \(z\) of the original-metric numerical range of
\(A_t^2-kA_t\) satisfies
\[
 \operatorname{Re}z\le\frac{E_N(t)^2-k^2}{4},\qquad
 (\operatorname{Im}z)^2\le E_N(t)^2
   \left(\frac{E_N(t)^2-k^2}{4}-\operatorname{Re}z\right).
\]
The proof retains every original input vector in homogeneous
coordinates and transports the complete inequality through
\(\Pi\) with metric \(\Pi^{-*}G_N\Pi^{-1}\).

The same source factors also control the calculated constituent
curvature. For its actual \(G_N\)-projection \(P_F\), set
\(a_\perp=\|(I-P_F)e_0\|_{G_N}^2\) and
\(b_F=\|P_FG_N^{-1}\ell^*\|_{G_N}^2\).
Let \(m(t),M(t)\) be the actual least and greatest eigenvalues of
\(G_N^{-1}\Pi^*\Pi\). Its positive curvature coefficient obeys
\[
 \frac{m(t)}{M(t)}\frac{a_\perp b_F}{|u|^2}
 \le c_F(t)\le
 \frac{M(t)}{m(t)}\frac{a_\perp b_F}{|u|^2}
 \le\frac{M(t)}{m(t)}\frac{\sigma_N^2}{|u|^2}.
\]
The complete proof supplies the exact nonnegative projection losses
between \(a_\perp b_F\) and \(\sigma_N^2\). These statements
retain the actual period-comparison eigenvalues and original theta
masses; no uniform packet-size bound is inserted.


\paragraph{Result R58: the original arithmetic quotient from a finite circle.}
The complete periodized-source proof (P1)--(P64), together with the
full source calculations (PSA1)--(PSA29) and (PSD1)--(PSD41), retains
$v_h=2\xi/h$, $w_h(t)=|v_h(1/2+it)|^2/(2\pi)$ and
$m_{h,k}=w_h^{*k}$. Its logarithmic coordinate map keeps the density
$e^{kr+\sum z_i}$ and the ordered-coordinate orientation
$(-1)^{k-1}$. The scalar circle sequence is related to the actual
relative Hilbert fibre by Fourier coefficients
$L^{-1/2}a_n\widehat{\mathcal U_kF_h^{\otimes k}}(2\pi n/L)$,
whose squared norm is $\sum_n\nu_n|a_n|^2$ with
$\nu_n=(2\pi/L)m_{h,k}(2\pi n/L)$. The zero mode remains in this sum.

For the original degree-$N$ polynomial Gram $M_N$ and unchanged
remainder map $J_N$, define the actual finite sampled Gram
$M_N(L,J)$ by retaining $|n|\le J$. The explicit original-source
quantities in (FC3)--(FC3a) produce
$\alpha=1-\eta_P-\eta_T>0$ and $\beta=1+\eta_P$, and prove
\[
 \alpha M_N\preceq M_N(L,J)\preceq\beta M_N,\qquad
 \alpha G_N\preceq G_N(L,J)\preceq\beta G_N,
 \quad G_N=(J_NM_N^{-1}J_N^*)^{-1}.
\]
The second assertion follows by taking minima over the same original
remainder fibres. The two minimizing representatives retain their
explicit difference in $\chi\mathcal P_{N-q}$ and its complete
theta primitive. On the common degree $D=2q+2$, all four endpoint
degrees and both generator raises are covered. Their exact error is
\[
 \left|\mathcal B_{h,k}(L,J)-\mathcal B_{h,k}\right|
      \le2q\log(\beta/\alpha).
\]
At fixed positive error budgets this is $O(q_k)$; its ratio to
$q_k\log k$ tends to zero. The required $L,J$ are explicit functions
of the original weighted and derivative Grams. Their growth and the
arithmetic sample enclosures are retained as actual quantities, rather
than replaced by unproved uniform estimates.

\paragraph{Result R59: the joint critical-jet and circle observation.}
The full density proof obtains the original exponential tails from
the gamma ray integral and the alternating Dirichlet series, including
every canceled value $v_h(\rho)=g^{(m)}(\rho)/h^{(m)}(\rho)\ne0$.
It proves density on the required analytic strip, with its
Cauchy--Schwarz half-width retained, and computes the completed circle
quotient exactly. If $b=p_L$ is the squarefree product of the sampled
roots of $\chi=\chi_{h,k}$, then
\[
 \ker\theta_L=(b)/(\chi),\qquad
 G_N(L)\downarrow G_\infty(L),\qquad
 \operatorname{rad}G_\infty(L)=(b)/(\chi).
\]
For $k\ge2$ the complete difference $G_N(L)-G_\infty(L)$ is
positive definite at every finite admitted degree: a nonzero
minimizing polynomial cannot vanish on the infinite complement of
the selected sample atoms. The limiting kernel retains the full
unsampled primary blocks and all nonconstant jets at sampled roots.

The proof (DC1)--(DC30) compares this with the actual weighted
tensor-critical map $C_k$, whose kernel is $(a)/(\chi)$ for
$a=\chi_{c,k}$. Put $d=\gcd(a,b)$ and $l=\operatorname{lcm}(a,b)$.
Using the complete critical inverse jets and the original weighted
circle coordinates gives
\[
 \operatorname{im}(C_k,\theta_L)
   \simeq\operatorname{im}C_k
        \mathop{\times}_{\mathbb C[S]/(d)}\operatorname{im}\theta_L
   \simeq\mathbb C[S]/(l).
\]
The joint kernel is $(l)/(\chi)$ and the common quotient is
$\mathbb C[S]/(d)$. At each original sum root, their quotient
lengths are respectively the maximum and minimum of the critical
cyclic length and the sampled length in $\{0,1\}$; all local unit
cofactors are retained in the kernel inclusions. A factorization
$\theta_L=T C_k$ exists exactly when $b\mid a$, with its explicit
inverse-unit and reduction formula. The reverse factorization exists
exactly when $a\mid b$. At $k=1$, full-order cancellation proves
that every sampled selected root has positive weight and $b\mid a$.
For general tensor degree, the extra circle summand has defining
polynomial $b/d$ and the explicit original polynomial source lift
in (DC29). Its existence criterion asserts no existence of an
off-critical zero.

All four observations transport through the same original period
matrix. For the role $s\in\{a,b,d,l\}$, retain its actual map $T_s$
and multiplier $M_s$. Then
\[
 T_{s,t}=T_s\Pi^{-1},\qquad
 u\partial_tT_{s,t}-M_sT_{s,t}
       =tT_s(e_0)\ell\Pi^{-1}.
\]
Their kernels, intersection, sum, crossed observations and complete
rank-one action defects are calculated in (DC23)--(DC30). Their
Split-Zero lifts retain each support and its receiving supported zero.

\paragraph{Result R60: universal residue endomorphisms and metric variation.}
The full calculation (RCX1)--(RCX40) extends the residue insertion
$R=e_0\ell$ over every commutative coefficient ring $B$ with monic
$\chi$. With the original residue matrix
$\mathcal B_{ij}=\ell(S^{i+j})$, its determinant remains
$(-1)^{q(q-1)/2}$. Every coefficient endomorphism has the exact inverse
expansion
\[
 T=\sum_{i,j}(T\mathcal B^{-1})_{ij}A^iRA^j.
\]
The full proof gives matrix units, their multiplication and base
change. Their original source lifts $r_NTj_E$ form an injective
representation in the corner with identity $r_Nj_E$ and kill the
same original theta boundaries.

For a proper nonzero invariant constituent $I:F\hookrightarrow E$,
retain the residual $r_M=(1-P_M)e_0$ and the residue-dual vector
$w_M=I(I^*MI)^{-1}(\ell I)^*$. They are $M$-orthogonal, with
squared norms $q_M,d_M>0$, and
$\mathfrak c_F(M)=q_Md_M$. Along any actual differentiable
positive-form path,
\[
 \frac{d}{ds}\log\mathfrak c_F(M)
 =\operatorname{Tr}\left[
   \left(\frac{r_Mr_M^*M}{q_M}-\frac{w_Mw_M^*M}{d_M}\right)
                  M^{-1}\dot M\right].
\]
The parenthesized operator has eigenvalues $1,-1,0$ with the
stated multiplicities, yielding the exact spectral-spread bound.
For the original endpoint forms $G_i\succeq G_j$, with extreme
relative eigenvalues $0<\theta\le\eta\le1$, this proves
\[
 \left|\log\frac{\mathfrak c_F(G_j)}{\mathfrak c_F(G_i)}\right|
 \le\log(\eta/\theta)\le-\log\theta
 \le\log\frac{\det G_i}{\det G_j}.
\]
The complete neighboring-minor formula
$\mathfrak c_F(M)=v_{J_+}(M)v_{J_-}(M)/v_{J_g}(M)^2$
retains every original frame determinant and both constituent
dimensions. The finite sampling comparison from R58 gives the
dimension-free coefficient interval
$\alpha/\beta\le\mathfrak c_F(G_N(L,J))/\mathfrak c_F(G_N)
\le\beta/\alpha$.

\paragraph{Result R61: sampled control with its exact commutator error.}
The full proof (FC1)--(FC29), including (FC3a), keeps the two actual
forms $G=G_N$, $H=G_N(L,J)$, the original
$T_t=A+tR-k/2$, and $B=G^{-1}(H-G)$. It first proves
\[
 (I+B)X_H=X_G+T_t^{\dagger_G}B+BT_t,
 \qquad X_G=T_t^{\dagger_G}+T_t,
 \quad X_H=H^{-1}(T_t^*H+HT_t).
\]
The coordinating Zeta calculation sharpens the transfer through
the explicit isometry $\mathcal S=(I+B)^{1/2}:(E,H)\to(E,G)$.
Set $Z=[\mathcal S,T_t]\mathcal S^{-1}$. Its complete Hermitian
error and the convergent Sylvester integral give
\[
 \left|\|X_H\|_H-\|X_G\|_G\right|
 \le\|Z+Z^{\dagger_G}\|_G
 \le\alpha^{-1}\|[B,A]+t[B,R]\|_G,
 \qquad [B,R]v=(Be_0)\ell(v)-e_0\ell(Bv).
\]
Both original metrics, the original unit and the ordered residue
terms remain. Commuting discrepancies give exactly equal control
radii. For the actual Fourier-projected representative
$\mathscr R$ and its theta boundary $\mathscr B$, the full identities
are
\[
 \begin{gathered}
 W_H(0)=-(\mathscr R^*\mathscr B+\mathscr B^*\mathscr R),\\
 (D_L^2-kD_L)\mathscr R-\mathscr R(A^2-kA)
       =(D_L-k)\mathscr B+\mathscr BA.
 \end{gathered}
\]
The resulting original-metric control radius gives both inequalities
of the full Laplacian numerical-range parabola (FC18), including the
zero-radius case. These statements preserve the actual arithmetic
commutator and boundary that a numerical endpoint estimate must
evaluate.

\paragraph{Result R62: quotient curvature and the complete fourth derivative.}
Retain the exact original coefficient sequence
$0\to F\xrightarrow I E\to E/F\to0$ and a constant quotient lift
$J_F$. Its least-period-norm lift and quotient Gram $K(t)$ are the
full Schur formulas (FC21). With $C=[I,J_F]$ they satisfy
\[
 \det(I^*\Pi^*\Pi I)\det K
      =|\det C|^2|\det\Pi|^2,
 \qquad
 \partial_{\bar t}\partial_t\log\det K
      =-\frac{|t|^2}{|u|^2}\mathfrak c_F(\Pi^*\Pi).
\]
The positive constituent curvature is therefore compensated in this
same coefficient quotient, with both curvatures vanishing at zero
and with opposite nonzero leading coefficients. The full proof
includes quotient-lift independence, basis covariance and the
literal comparison with the original theta quotient metric.

The residue continuation evaluates the pure jets as well. If
$\psi=\log\det(I^*\Pi(t)^*\Pi(t)I)$ and
$\phi(t)=\log\det(I^*\Pi(0)^*\Pi(t)I)$ near zero, then
\[
 \left.\frac{d^4}{dx^4}\psi(x)\right|_0
   =2\operatorname{Re}\phi^{(4)}(0)
       +6\mathfrak c_F(\Pi(0)^*\Pi(0))/|u|^2.
\]
The full ordered recursion for $\Pi^{(n)}=\Pi Q_n$, its explicit
$Q_4$ including $3R^2/u^2$, and the complete trace-log formula in
(RCX30)--(RCX32) evaluate the remaining pure term. The forced
period equation retains $u\Pi R$, and its original source
primitive retains $tKR$. The exact Pascal translation preserves
these derivatives of order at least two while its scalar factor
retains the first-derivative trace shift $2p\operatorname{Re}a$.


The arithmetic norm theorem, sharp endpoint products and upper
comparison therefore enter the existing Toda, Gamma and exterior
maps with complete proof texts. The finite upper determinants in
R45 are the concrete next quantities to estimate at the amplified
degree; the explicit scalar estimate and the lower limit $4$
currently coexist with the gap proved above.
